\documentclass[12pt,a4paper,oneside]{extreport}
\usepackage{ctu-report}
\usepackage{amsmath}
\usepackage{amssymb}
\usepackage{placeins}
\usepackage{needspace}
\usetikzlibrary{positioning}
\setlength{\emergencystretch}{3em}
\renewcommand{\lstlistingname}{Đoạn mã}
\AtBeginDocument{\renewcommand{\thelstlisting}{\arabic{chapter}.\arabic{lstlisting}}}
\graphicspath{{assets/figures/}}
\Urlmuskip=0mu plus 1mu
\def\UrlBreaks{\do\/\do-\do.\do_}

% Figure helper: \ctufigure{file}{width as a fraction of \textwidth}{caption}{label}
% Height is capped so tall screenshots still fit on a page with their caption.
\newcommand{\ctufigure}[4]{%
  \begin{figure}[!htbp]
    \centering
    \includegraphics[width=#2\textwidth,height=0.62\textheight,keepaspectratio]{#1}
    \caption{#3}
    \label{#4}
  \end{figure}
  \FloatBarrier
}

% Caption for the unnumbered tables in the front matter and introduction.
\newcommand{\ctuplaintablecaption}[1]{%
  \begin{center}
    \itshape Bảng: #1
  \end{center}
  \addcontentsline{lot}{table}{Bảng: #1}%
}

% Khung tiểu luận: hộp TODO đánh dấu phần cần viết. Xóa macro này khi hoàn thiện.
\newcommand{\ctutodo}[1]{%
  \par\smallskip\noindent
  \fcolorbox{red!60!black}{red!5}{\parbox{\dimexpr\linewidth-2\fboxsep-2\fboxrule\relax}{%
    \small\color{red!60!black}\textbf{TODO:} #1}}%
  \par\smallskip
}

% -----------------------------------------------------------------------------
% METADATA
% -----------------------------------------------------------------------------
\CTUReportType{TIỂU LUẬN}
\CTUTitle{Xây dựng quy trình DevSecOps an toàn với cơ chế triển khai không gián đoạn cho ứng dụng Web Java}
\CTUStudent{Trần Thanh Duy}
\CTUStudentID{DC24V7X604}
\CTUCohort{2024--2027}
\CTUSupervisor{ThS. Nguyễn Minh Trung} % TODO: xác nhận GVHD tiểu luận
\CTUMonthYear{MM/YYYY} % TODO: tháng/năm nộp

\begin{document}

\ctumakecover
\ctumakeinnertitle

% -----------------------------------------------------------------------------
% FRONT MATTER
% -----------------------------------------------------------------------------
\ctufrontmatter

\ctuheading{LỜI CẢM ƠN}
Em xin gửi lời cảm ơn chân thành đến các thầy cô tại Đại học Cần Thơ, đặc biệt
là các thầy cô trong trường Công Nghệ Thông Tin và Truyền Thông. Trong thời
gian qua các thầy cô đã tận tình giảng dạy, hướng dẫn nhiệt tình để em có được
những kiến thức mới đầy hữu ích cho tương lai sau này.

Em xin cảm ơn Thầy Nguyễn Minh Trung đã hỗ trợ hướng dẫn em hoàn thành tiểu
luận ``Xây dựng quy trình DevSecOps an toàn với cơ chế triển khai không gián
đoạn cho ứng dụng Web Java''.

Do trình độ lý luận và kinh nghiệm thực tế còn hạn chế nên bài báo cáo còn
nhiều thiếu sót, em rất mong nhận được ý kiến đóng góp từ Thầy, Cô để em có thể
rút ra nhiều kinh nghiệm hơn.

Cuối cùng, em xin chân thành cảm ơn gia đình và bạn bè, đã luôn tạo điều kiện,
quan tâm, giúp đỡ, động viên em trong suốt quá trình học tập và hoàn thành bài
tiểu luận này.

\begin{flushright}
  Em xin chân thành cảm ơn
\end{flushright}

\ctutableofcontents
\ctulistoffigures
\ctulistoftables

\ctuheading{DANH MỤC TỪ VIẾT TẮT}
\begin{longtable}{|>{\raggedright\arraybackslash}p{0.24\textwidth}|>{\raggedright\arraybackslash}p{0.68\textwidth}|}
  \hline
  \ctuTableHead{Kí hiệu \& Chữ viết tắt} & \ctuTableHead{Diễn giải} \\
  \hline
  \endhead
  API & Application Programming Interface (Giao diện lập trình ứng dụng) \\
  \hline
  CD & Continuous Delivery / Continuous Deployment (Phân phối liên tục / Triển khai liên tục) \\
  \hline
  CI & Continuous Integration (Tích hợp liên tục) \\
  \hline
  CVE & Common Vulnerabilities and Exposures (Danh mục lỗ hổng bảo mật công khai) \\
  \hline
  CVSS & Common Vulnerability Scoring System (Hệ thống chấm điểm mức độ nghiêm trọng lỗ hổng) \\
  \hline
  DevSecOps & Development, Security and Operations (Tích hợp bảo mật vào DevOps) \\
  \hline
  DORA & DevOps Research and Assessment (Bộ chỉ số đo hiệu quả phân phối phần mềm) \\
  \hline
  GHCR & GitHub Container Registry (Kho lưu trữ Container của GitHub) \\
  \hline
  JDK & Java Development Kit (Bộ công cụ phát triển phần mềm Java) \\
  \hline
  MTTR & Mean Time To Recovery (Thời gian khôi phục trung bình) \\
  \hline
  OCI & Open Container Initiative (Tiêu chuẩn mở cho định dạng container) \\
  \hline
  OWASP & Open Worldwide Application Security Project \\
  \hline
  SAST & Static Application Security Testing (Kiểm thử bảo mật tĩnh) \\
  \hline
  SBOM & Software Bill of Materials (Danh mục thành phần phần mềm) \\
  \hline
  SCA & Software Composition Analysis (Phân tích thành phần phần mềm) \\
  \hline
  SLSA & Supply-chain Levels for Software Artifacts \\
  \hline
  SSDF & Secure Software Development Framework (NIST SP 800-218) \\
  \hline
  STRIDE & Spoofing, Tampering, Repudiation, Information disclosure, Denial of service, Elevation of privilege \\
  \hline
  vCPU & Virtual Central Processing Unit (Nhân CPU ảo) \\
  \hline
  YAML & YAML Ain't Markup Language (Ngôn ngữ cấu hình dữ liệu) \\
  \hline
\end{longtable}
\ctuplaintablecaption{Danh mục từ viết tắt}

% -----------------------------------------------------------------------------
% MAIN MATTER
% -----------------------------------------------------------------------------
\ctumainmatter

\ctuheading{TÓM TẮT}
\ctutodo{Viết sau cùng (khoảng 250--300 từ), theo cấu trúc: bối cảnh (CI/CD phổ
biến nhưng thường bỏ qua bảo mật chuỗi cung ứng và tính sẵn sàng khi triển
khai) $\rightarrow$ vấn đề (pipeline niên luận: cổng bảo mật không chặn, Watchtower
gây gián đoạn và không rollback) $\rightarrow$ giải pháp (các cổng SAST/SCA/quét
image/ký số + Blue-Green có health check và tự rollback) $\rightarrow$ phương pháp
đánh giá (4 thực nghiệm E1--E4) $\rightarrow$ kết quả chính bằng số liệu.}

\textbf{Từ khóa:} DevSecOps, CI/CD, bảo mật chuỗi cung ứng phần mềm, SBOM, ký
số container image, Blue-Green Deployment, GitHub Actions.

\ctuheading{GIỚI THIỆU}
\introsection{Đặt vấn đề}
Trong bối cảnh ngành công nghiệp phần mềm hiện đại, phương pháp phát triển phần
mềm Agile đang trở thành tiêu chuẩn cốt lõi tại các công ty công nghệ lớn nhằm
đáp ứng nhu cầu thay đổi nhanh chóng của thị trường. Thông qua việc chia nhỏ quá
trình phát triển thành các chu kỳ ngắn (Sprints), đội ngũ lập trình
(Developer/Dev) có thể liên tục chuyển giao các tính năng phần mềm mới một cách
nhanh chóng.

Tuy nhiên, tốc độ lập trình và tần suất chuyển giao gia tăng này tạo ra một nút
thắt lớn ở giai đoạn triển khai (Deployment). Nó đặt áp lực nặng nề lên đội ngũ
vận hành (Operations/Ops). Ngoài việc phải đáp ứng khối lượng lớn các yêu cầu
cập nhật liên tục từ đội Dev, team Ops còn phải chịu trách nhiệm đảm bảo thời
gian hoạt động (uptime), tính sẵn sàng và sự ổn định của hệ thống máy chủ tổng
thể.

Sự mâu thuẫn giữa tốc độ phát triển và sự ổn định trong vận hành bộc lộ rõ qua
những điểm yếu của quy trình làm việc thủ công truyền thống:
\begin{itemize}
  \item \textbf{Ở giai đoạn Phát Triển (Development):} Thông thường, các nhóm
  phát triển làm việc độc lập trong nhiều tuần với khối lượng mã nguồn (source
  code) đồ sộ. Khi toàn bộ mã nguồn này được gộp chung và đưa lên máy chủ để
  biên dịch (build) cùng một lúc, hệ thống thường xuyên phát sinh hàng loạt lỗi
  phần mềm và xung đột cấu hình. Việc phát hiện lỗi quá muộn ở giai đoạn cuối
  khiến các kỹ sư mất rất nhiều thời gian để rà soát và khắc phục sự cố.
  \item \textbf{Ở giai đoạn Triển khai (Deployment):} Việc đưa ứng dụng lên môi
  trường thực tế (Production) là một chuỗi các công việc phức tạp, bao gồm cung
  cấp máy chủ, cấu hình các thư viện phụ thuộc (dependencies), thay đổi thiết
  lập mạng và triển khai các gói phần mềm đóng gói (artifact như định dạng
  jar/war, zip/tar). Khi nhóm Ops liên tục nhận các yêu cầu triển khai thủ công
  nhưng lại thiếu đi tài liệu hướng dẫn cấu hình rõ ràng, rủi ro về sai sót của
  con người là rất lớn, dẫn đến nguy cơ gián đoạn dịch vụ và giảm trải nghiệm
  người dùng.
\end{itemize}

Để giải quyết triệt để sự mâu thuẫn này, văn hóa DevOps ra đời nhằm phá bỏ rào
cản truyền thống, giúp đội ngũ Development và Operations làm việc gắn kết chặt
chẽ với nhau. Trọng tâm của DevOps là áp dụng các công cụ tự động hóa
(Automation Tools) nhằm giảm thiểu thao tác thủ công và loại bỏ sai sót từ con
người. Cốt lõi của mô hình này là quy trình Tích hợp liên tục và Phân
phối/triển khai liên tục (CI/CD):
\begin{itemize}
  \item Tích hợp liên tục (Continuous Integration -- CI): Yêu cầu mọi thay đổi
  (commit/push) mã nguồn của lập trình viên phải đi kèm với quá trình xây dựng
  và kiểm thử tự động ngay lập tức. Việc này giúp phát hiện và cô lập lỗi (sự
  cố nền, xung đột) ngay từ giai đoạn đầu để chúng không nhân lên diện rộng. Sau
  khi vượt qua kiểm định, mã nguồn được đóng gói thành phần mềm hoàn chỉnh
  (artifact). CI hoạt động như một hệ thống kiểm tra tự động. Mỗi khi code được
  thay đổi và commit lên hệ thống, các bài kiểm tra tự động sẽ được chạy để
  kiểm tra xem code có hoạt động đúng và không gây lỗi hay không.
  \item Phân phối liên tục (Continuous Delivery): là quá trình tiếp nối CI,
  trong đó mọi phiên bản phần mềm đã vượt qua các bước build, kiểm thử và đánh
  giá bảo mật sẽ được tự động đóng gói, lưu trữ tại kho artifact hoặc container
  registry, đồng thời sẵn sàng để triển khai lên các môi trường như Testing,
  Staging hoặc Production. Với Continuous Delivery, việc triển khai lên
  Production có thể cần sự phê duyệt thủ công của người quản trị.
  \item Triển khai liên tục (Continuous Deployment): là mức độ tự động hóa cao
  hơn của Continuous Delivery. Sau khi mã nguồn vượt qua toàn bộ các bước kiểm
  định trong pipeline, phiên bản mới sẽ được tự động triển khai lên môi trường
  chạy thực tế mà không cần thao tác thủ công. Để bảo đảm an toàn, quy trình
  này cần kết hợp cơ chế kiểm tra tình trạng hoạt động của ứng dụng (health
  check), theo dõi nhật ký hệ thống, kiểm soát phiên bản image và khả năng quay
  lui về phiên bản ổn định trước đó khi xảy ra sự cố.
\end{itemize}


Trong niên luận trước, tác giả đã xây dựng một quy trình CI/CD cho ứng dụng Web
Java với GitHub Actions, GitHub Container Registry và Watchtower. Tuy nhiên, quá
trình vận hành và đánh giá cho thấy hai nhóm hạn chế:
\begin{itemize}
  \item \textbf{Về bảo mật:} các bước kiểm tra chỉ sinh báo cáo mà không chặn
  pipeline (Trivy cấu hình \texttt{exit-code: 0}, Checkstyle không làm build
  thất bại), chưa có phân tích mã tĩnh, chưa có SBOM và image không được ký nên
  máy chủ không thể xác minh nguồn gốc image trước khi chạy.
  \item \textbf{Về triển khai:} Watchtower dừng container cũ rồi mới khởi động
  container mới nên ứng dụng bị gián đoạn trong mỗi lần cập nhật; Watchtower
  không kiểm tra tình trạng ứng dụng sau khi khởi động và không có cơ chế quay
  lui khi phiên bản mới bị lỗi; việc theo dõi tag \texttt{latest} khiến khó
  truy vết phiên bản đang chạy.
\end{itemize}

\ctutodo{Bổ sung 1--2 đoạn dẫn chứng bối cảnh: các sự cố chuỗi cung ứng phần
mềm tiêu biểu (SolarWinds 2020, Log4Shell 2021, xz-utils 2024) và số liệu từ
báo cáo DORA về ảnh hưởng của tỉ lệ triển khai lỗi; trích dẫn nguồn.}

Xuất phát từ các hạn chế trên, tiểu luận ``Xây dựng quy trình DevSecOps an toàn
với cơ chế triển khai không gián đoạn cho ứng dụng Web Java'' được thực hiện
nhằm nâng cấp pipeline của niên luận theo hai hướng: tích hợp các cổng kiểm soát
bảo mật xuyên suốt pipeline và thay thế cơ chế cập nhật container bằng chiến
lược Blue-Green có kiểm tra sức khỏe và tự động quay lui.

\introsection{Mục tiêu đề tài}
\introsubsection{Mục tiêu tổng quát}
Xây dựng và đánh giá thực nghiệm một quy trình DevSecOps cho ứng dụng Web Java,
trong đó mọi phiên bản phải vượt qua các cổng kiểm soát bảo mật, được ký số và
xác minh trước khi triển khai, và được đưa vào vận hành mà không gây gián đoạn
dịch vụ.

\introsubsection{Câu hỏi nghiên cứu}
\begin{itemize}
  \item \textbf{RQ1:} Việc bổ sung các cổng kiểm soát bảo mật (SAST, SCA, quét
  secret, quét container image, ký số và xác minh image) ngăn chặn được bao
  nhiêu loại lỗ hổng trước khi chúng đến môi trường vận hành, và làm tăng thời
  gian thực thi pipeline bao nhiêu?
  \item \textbf{RQ2:} So với cơ chế cập nhật của Watchtower, chiến lược
  Blue-Green kết hợp health check và tự động rollback cải thiện thời gian gián
  đoạn dịch vụ và khả năng khôi phục (MTTR) đến mức nào?
\end{itemize}

\introsubsection{Đóng góp của đề tài}
\begin{itemize}
  \item Phân tích rủi ro pipeline CI/CD hiện tại theo mô hình STRIDE và ánh xạ
  với các khuyến nghị của NIST SSDF, SLSA và OWASP Top 10 CI/CD Security Risks.
  \item Thiết kế và cài đặt pipeline DevSecOps với ba lớp cổng kiểm soát, sinh
  SBOM và ký số image bằng Cosign.
  \item Thiết kế và cài đặt cơ chế triển khai Blue-Green trên một máy chủ đơn
  với Docker Compose và Nginx, có kiểm tra sức khỏe và tự động quay lui.
  \item Bộ thực nghiệm có thể lặp lại (E1--E4) cùng số liệu so sánh giữa
  pipeline cũ và pipeline mới.
\end{itemize}

\introsection{Đối tượng và phạm vi nghiên cứu}
\introsubsection{Đối tượng nghiên cứu}
Quy trình DevSecOps và bảo mật chuỗi cung ứng phần mềm trong CI/CD; các chiến
lược triển khai ứng dụng (Recreate, Rolling, Blue-Green, Canary); các công cụ
GitHub Actions, Docker, GHCR, Trivy, OWASP Dependency-Check, CodeQL, Gitleaks,
Syft, Cosign, Nginx và k6.

\introsubsection{Phạm vi nghiên cứu}
\begin{itemize}
  \item Ứng dụng mẫu VProfile (Java Spring MVC) được giữ nguyên nghiệp vụ; đề
  tài chỉ bổ sung kiểm thử tự động cần thiết cho pipeline.
  \item Môi trường triển khai là một máy ảo Ubuntu 24.04 tạo bằng Vagrant, chạy
  Docker Compose. Đề tài không sử dụng Kubernetes và không triển khai lên hạ
  tầng đám mây.
  \item Cơ sở dữ liệu dùng chung giữa hai môi trường Blue và Green; đề tài không
  giải quyết bài toán migration lược đồ cơ sở dữ liệu không tương thích ngược.
\end{itemize}

\introsection{Phương pháp nghiên cứu}
\begin{itemize}
  \item \textbf{Nghiên cứu tài liệu:} tổng hợp các nghiên cứu và tiêu chuẩn về
  DevSecOps, bảo mật chuỗi cung ứng và chiến lược triển khai.
  \item \textbf{Thiết kế -- cài đặt:} mở rộng pipeline của niên luận theo thiết
  kế ở Chương~\ref{chap:design}.
  \item \textbf{Thực nghiệm so sánh:} chạy pipeline cũ (baseline) và pipeline mới
  trên cùng điều kiện, lặp lại nhiều lần và báo cáo giá trị trung bình, độ lệch
  chuẩn.
\end{itemize}

\introsection{Công nghệ sử dụng}
\begin{center}
  \begin{tabular}{|>{\raggedright\arraybackslash}p{0.40\textwidth}|>{\raggedright\arraybackslash}p{0.50\textwidth}|}
    \hline
    \ctuTableHead{Nhóm công nghệ} & \ctuTableHead{Công nghệ} \\
    \hline
    Quản lý mã nguồn, CI/CD & Git, GitHub, GitHub Actions \\
    \hline
    Build, kiểm thử, độ phủ & Maven, JUnit, JaCoCo, Checkstyle \\
    \hline
    Phân tích mã tĩnh (SAST) & CodeQL / SpotBugs + FindSecBugs \\
    \hline
    Phân tích thư viện (SCA) & OWASP Dependency-Check \\
    \hline
    Quét secret & Gitleaks \\
    \hline
    Quét container image & Trivy \\
    \hline
    SBOM, ký số & Syft, Cosign (Sigstore) \\
    \hline
    Container, registry & Docker, Docker Compose, GHCR \\
    \hline
    Triển khai & Nginx (Blue-Green), GitHub Actions self-hosted runner, Watchtower (baseline) \\
    \hline
    Môi trường & Ubuntu 24.04 trên Vagrant (VMware/VirtualBox) \\
    \hline
    Đo lường & k6 \\
    \hline
  \end{tabular}
\end{center}
\ctuplaintablecaption{Các công nghệ được sử dụng}
\ctutodo{Chốt lựa chọn CodeQL hay SpotBugs + FindSecBugs sau khi thử trên mã
nguồn VProfile, cập nhật lại bảng.}

\introsection{Bố cục của tiểu luận}
Tiểu luận được chia làm 5 chương:
\begin{itemize}
  \item Chương I: Tổng quan và các nghiên cứu liên quan
  \item Chương II: Cơ sở lý thuyết
  \item Chương III: Phân tích và thiết kế hệ thống
  \item Chương IV: Cài đặt hệ thống
  \item Chương V: Thực nghiệm và đánh giá
\end{itemize}
Phần Kết luận tổng hợp câu trả lời cho các câu hỏi nghiên cứu, nêu hạn chế và
hướng phát triển.

% -----------------------------------------------------------------------------
% CHƯƠNG I
% -----------------------------------------------------------------------------
\chapter{TỔNG QUAN VÀ CÁC NGHIÊN CỨU LIÊN QUAN}
\label{chap:related}
Chương này trình bày bối cảnh của đề tài thông qua ba nhóm tài liệu: các sự cố
và nghiên cứu về bảo mật chuỗi cung ứng phần mềm, các nghiên cứu về DevSecOps,
và các nghiên cứu về chiến lược triển khai phần mềm. Tiếp theo, chương giới
thiệu các tiêu chuẩn được dùng làm khung tham chiếu cho thiết kế ở
Chương~\ref{chap:design}. Cuối chương là phần so sánh các công trình liên quan
và xác định khoảng trống mà tiểu luận hướng tới.

% -----------------------------------------------------------------------------
\section{Bảo mật chuỗi cung ứng phần mềm}
\label{sec:supply-chain}
\subsection{Khái niệm}
Chuỗi cung ứng phần mềm (software supply chain) là tập hợp toàn bộ các thành
phần, con người, quy trình và hạ tầng tham gia vào việc tạo ra và phân phối một
sản phẩm phần mềm: mã nguồn do nhóm phát triển viết, các thư viện mã nguồn mở
của bên thứ ba, công cụ và máy chủ build, kho lưu trữ artifact hoặc container
registry, cho đến hạ tầng triển khai và vận hành.

Ứng dụng VProfile trong đề tài là một ví dụ điển hình: phần mã nghiệp vụ do nhóm
phát triển viết chỉ chiếm một phần nhỏ, phần lớn mã thực thi đến từ các thư viện
như Spring Framework, Spring Security, Hibernate, các trình kết nối MySQL,
RabbitMQ, Memcached và từ image nền Tomcat. Mỗi thành phần này lại có các phụ
thuộc bắc cầu (transitive dependencies) của riêng nó. Do đó, kẻ tấn công không
nhất thiết phải khai thác trực tiếp mã nguồn của ứng dụng mà có thể xâm nhập ở
bất kỳ mắt xích nào trong chuỗi (Hình~\ref{fig:supply-chain}).

\begin{figure}[!htbp]
  \centering
  \begin{tikzpicture}[
      node distance=4mm,
      link/.style={draw, rounded corners, align=center, minimum height=12mm,
        text width=22mm, font=\small},
      inc/.style={align=center, text width=22mm, font=\footnotesize\itshape,
        text=red!60!black},
      >=stealth]
    \node[link] (src) {Mã nguồn};
    \node[link, right=of src] (dep) {Thư viện phụ thuộc};
    \node[link, right=of dep] (build) {Hệ thống build};
    \node[link, right=of build] (art) {Artifact / Registry};
    \node[link, right=of art] (dep2) {Triển khai};
    \foreach \a/\b in {src/dep, dep/build, build/art, art/dep2}
      \draw[->, thick] (\a) -- (\b);
    \node[inc, below=3mm of src] {Lộ secret, mã độc do người đóng góp};
    \node[inc, below=3mm of dep] {Log4Shell, xz-utils};
    \node[inc, below=3mm of build] {SolarWinds};
    \node[inc, below=3mm of art] {Image bị thay thế, không xác minh};
    \node[inc, below=3mm of dep2] {Triển khai bản lỗi, không quay lui};
  \end{tikzpicture}
  \caption{Các mắt xích của chuỗi cung ứng phần mềm và ví dụ rủi ro tương ứng}
  \label{fig:supply-chain}
\end{figure}
\FloatBarrier

\subsection{Các sự cố tiêu biểu}
\textbf{SolarWinds (2020).} Nhóm tấn công đã xâm nhập vào hệ thống build của
phần mềm giám sát mạng SolarWinds Orion và chèn mã độc (được đặt tên SUNBURST)
vào trong quá trình biên dịch~\cite{mandiant-sunburst}. Vì mã độc được thêm vào
ở giai đoạn build, các bản cập nhật chứa mã độc vẫn được ký số hợp lệ bởi nhà
cung cấp và được phân phối qua kênh cập nhật chính thức. Theo công bố của
SolarWinds, số khách hàng đã cài đặt bản cập nhật bị nhiễm vào khoảng dưới
18.000 tổ chức. Sự cố cho thấy chữ ký số trên artifact là chưa đủ nếu bản thân
hệ thống build không đáng tin cậy, đây là động lực trực tiếp cho khung SLSA
(Mục~\ref{sec:standards}).

\textbf{Log4Shell (CVE-2021-44228, 2021).} Lỗ hổng thực thi mã từ xa trong thư
viện ghi log Apache Log4j~2, khai thác thông qua cơ chế tra cứu JNDI trong chuỗi
được ghi log, có điểm CVSS tối đa 10.0~\cite{nvd-log4shell}. Mức độ ảnh hưởng
của Log4Shell rất rộng vì Log4j thường xuất hiện như một phụ thuộc bắc cầu:
nhiều tổ chức không biết hệ thống của mình có sử dụng thư viện này hay không.
Sự cố nhấn mạnh vai trò của phân tích thành phần phần mềm (SCA) và danh mục
thành phần phần mềm (SBOM) để trả lời nhanh câu hỏi ``hệ thống nào đang bị ảnh
hưởng''.

\textbf{xz-utils (CVE-2024-3094, 2024).} Một tài khoản đã đóng góp cho dự án
nén dữ liệu xz-utils trong thời gian dài, trở thành người duy trì (maintainer)
được tin tưởng, sau đó chèn cửa hậu vào các gói phát hành phiên bản 5.6.0 và
5.6.1~\cite{nvd-xz}. Cửa hậu nhắm vào tiến trình \texttt{sshd} trên một số bản
phân phối Linux thông qua thư viện \texttt{liblzma}. Mã độc chỉ có trong gói
phát hành, không có trong mã nguồn hiển thị trên kho Git, và được phát hiện tình
cờ khi một kỹ sư điều tra độ trễ bất thường của SSH. Sự cố cho thấy rủi ro đến
từ chính các dự án thượng nguồn (upstream) và tầm quan trọng của việc quét cả
artifact cuối cùng, không chỉ mã nguồn.

\subsection{Phản ứng của cộng đồng và chính phủ}
Sau các sự cố trên, bảo mật chuỗi cung ứng trở thành ưu tiên ở cấp chính sách.
Sắc lệnh hành pháp 14028 của Hoa Kỳ (05/2021) yêu cầu nhà cung cấp phần mềm cho
chính phủ liên bang tuân thủ các thực hành phát triển an toàn và cung cấp
SBOM~\cite{eo14028}. Viện Tiêu chuẩn và Công nghệ Quốc gia Hoa Kỳ (NIST) ban hành
SSDF~\cite{nist-ssdf}; Quỹ OpenSSF phát triển khung SLSA~\cite{slsa} và dự án
Sigstore~\cite{sigstore-paper} nhằm giúp việc ký số artifact trở nên phổ biến.
Bảng~\ref{tab:incidents} tóm tắt các sự cố, mắt xích bị khai thác và biện pháp
tương ứng được áp dụng trong đề tài.

\begin{table}[!htbp]
  \centering
  \fontsize{11pt}{13.42pt}\selectfont
  \renewcommand{\arraystretch}{1.25}
  \begin{tabularx}{\textwidth}{|>{\raggedright\arraybackslash}p{0.17\textwidth}|>{\raggedright\arraybackslash}p{0.17\textwidth}|>{\raggedright\arraybackslash}X|>{\raggedright\arraybackslash}X|}
    \hline
    \ctuTableHead{Sự cố} & \ctuTableHead{Mắt xích} & \ctuTableHead{Bài học} & \ctuTableHead{Biện pháp trong đề tài} \\
    \hline
    SolarWinds & Hệ thống build & Chữ ký không đủ nếu build bị xâm nhập & Build trên runner do GitHub quản lý, ký keyless gắn với danh tính workflow \\
    \hline
    Log4Shell & Thư viện phụ thuộc & Cần biết chính xác hệ thống dùng thư viện nào & SCA (OWASP Dependency-Check), SBOM (Syft) \\
    \hline
    xz-utils & Thư viện / gói phát hành & Mã độc có thể chỉ nằm trong artifact & Quét container image (Trivy) sau khi build \\
    \hline
    Image bị thay thế & Registry & Máy chủ cần xác minh nguồn gốc image & Ký (Cosign) và \texttt{cosign verify} trước khi triển khai \\
    \hline
  \end{tabularx}
  \caption{Các sự cố chuỗi cung ứng tiêu biểu và biện pháp tương ứng}
  \label{tab:incidents}
\end{table}
\FloatBarrier

% -----------------------------------------------------------------------------
\section{Các nghiên cứu về DevSecOps}
\label{sec:devsecops-research}
\subsection{Định nghĩa và mức độ trưởng thành của DevSecOps}
Mohan và Ben Othmane~\cite{mohan2016} là một trong những công trình sớm khảo sát
có hệ thống về việc tích hợp bảo mật vào DevOps (SecDevOps). Bằng phương pháp
mapping study, các tác giả nhận thấy lĩnh vực này khi đó còn non trẻ: nhiều tài
liệu đến từ giới thực hành (blog, bài thuyết trình) hơn là nghiên cứu học
thuật, và chưa có định nghĩa thống nhất.

Myrbakken và Colomo-Palacios~\cite{myrbakken2017} thực hiện một tổng quan đa
nguồn (multivocal literature review), kết hợp cả tài liệu học thuật và tài liệu
thực hành, để làm rõ định nghĩa, đặc điểm, lợi ích và thách thức của DevSecOps.
Nghiên cứu chỉ ra DevSecOps không chỉ là bổ sung công cụ bảo mật vào pipeline
mà còn là thay đổi văn hóa, trong đó trách nhiệm bảo mật được chia sẻ giữa
nhóm phát triển, vận hành và bảo mật; tự động hóa các bước kiểm tra bảo mật là
một trong các đặc điểm cốt lõi.

\subsection{Thách thức khi áp dụng DevSecOps}
Rajapakse và cộng sự~\cite{rajapakse2022} thực hiện tổng quan có hệ thống (SLR)
các công trình đã được bình duyệt về thách thức và giải pháp khi áp dụng
DevSecOps. Các thách thức được nhóm theo bốn chủ đề: con người, thực hành, công
cụ và hạ tầng. Đáng chú ý với đề tài này là nhóm thách thức về công cụ: các
công cụ kiểm thử bảo mật tự động dễ sinh cảnh báo sai (false positive), làm
chậm pipeline, và khó tích hợp vào luồng CI/CD vốn yêu cầu phản hồi nhanh. Đây
chính là sự đánh đổi mà thực nghiệm E1 và E2 của tiểu luận đo lường.

Enck và Williams~\cite{enck2022} tổng hợp kết quả từ các buổi thảo luận với
nhiều tổ chức trong công nghiệp và chính phủ, đưa ra năm thách thức hàng đầu
của bảo mật chuỗi cung ứng: cập nhật các phụ thuộc có lỗ hổng, khai thác SBOM
cho mục đích bảo mật, lựa chọn thành phần phụ thuộc đáng tin cậy, bảo vệ quy
trình build, và mở rộng việc áp dụng các biện pháp trên ở quy mô lớn.

\subsection{Tấn công và phòng vệ chuỗi cung ứng}
Ohm và cộng sự~\cite{ohm2020} thu thập một tập dữ liệu các gói mã độc thực tế
trên các hệ sinh thái npm, PyPI và RubyGems, phân tích cách chúng được đưa vào
(ví dụ đặt tên gần giống gói phổ biến -- typosquatting, chiếm quyền tài khoản
người duy trì) và thời điểm mã độc được kích hoạt. Ladisa và
cộng sự~\cite{ladisa2023} đề xuất một phân loại (taxonomy) các hướng tấn công
vào chuỗi cung ứng mã nguồn mở dưới dạng cây tấn công, liên kết chúng với các
sự cố thực tế và các biện pháp bảo vệ, đồng thời đánh giá các biện pháp này qua
khảo sát chuyên gia.

Về phía phòng vệ, Newman và cộng sự~\cite{sigstore-paper} giới thiệu Sigstore,
hệ thống ký số ``không cần khóa'' (keyless): thay vì quản lý khóa bí mật dài
hạn, người ký xác thực bằng danh tính OpenID Connect (ví dụ danh tính của một
workflow GitHub Actions), nhận chứng chỉ ngắn hạn, và mọi chữ ký được ghi vào
một nhật ký minh bạch (transparency log) công khai. Cách tiếp cận này loại bỏ
rủi ro lộ khóa ký và là cơ sở cho việc ký image bằng Cosign trong đề tài.

% -----------------------------------------------------------------------------
\section{Các nghiên cứu về chiến lược triển khai}
\label{sec:deploy-research}
Shahin và cộng sự~\cite{shahin2017} thực hiện tổng quan có hệ thống về các cách
tiếp cận, công cụ, thách thức và thực hành của tích hợp, phân phối và triển
khai liên tục. Nghiên cứu ghi nhận rằng triển khai liên tục (Continuous
Deployment) khó áp dụng hơn đáng kể so với tích hợp liên tục, một phần do rủi
ro khi đưa phiên bản mới vào môi trường vận hành mà không có sự kiểm soát của
con người; do đó các cơ chế giảm thiểu rủi ro khi triển khai là một hướng cần
được quan tâm.

Humble và Farley~\cite{humble2010} hệ thống hóa các mẫu triển khai giúp giảm
rủi ro khi phát hành, trong đó có Blue-Green Deployment (duy trì hai môi trường
giống hệt nhau, chuyển lưu lượng từ môi trường cũ sang môi trường mới chỉ khi
môi trường mới đã sẵn sàng) và Canary Release (đưa phiên bản mới đến một phần
nhỏ người dùng trước). Fowler~\cite{fowler-bluegreen} nhấn mạnh ưu điểm then
chốt của Blue-Green là khả năng quay lui gần như tức thời: chỉ cần chuyển lưu
lượng trở lại môi trường cũ vẫn đang chạy.

Forsgren, Humble và Kim~\cite{forsgren2018}, dựa trên dữ liệu khảo sát nhiều
năm của chương trình nghiên cứu DORA~\cite{dora}, đề xuất bốn chỉ số then chốt
để đo hiệu quả phân phối phần mềm: tần suất triển khai, thời gian từ khi thay
đổi mã đến khi vận hành, tỉ lệ thay đổi gây lỗi và thời gian khôi phục dịch vụ.
Một phát hiện quan trọng của chương trình là tốc độ và độ ổn định không phải là
hai mục tiêu đánh đổi lẫn nhau: các tổ chức hiệu quả nhất đạt được đồng thời cả
hai. Kết quả này ủng hộ giả thuyết của đề tài rằng việc tăng tự động hóa (cổng
bảo mật, health check, rollback) có thể cải thiện độ an toàn mà không làm mất
đi lợi ích của triển khai liên tục.

Trong thực tế công nghiệp, các chiến lược Blue-Green và Canary thường được hiện
thực trên nền tảng điều phối container như Kubernetes cùng các công cụ như Argo
Rollouts. Ngược lại, với hạ tầng nhỏ chỉ gồm một máy chủ chạy Docker Compose
--- vốn phổ biến ở doanh nghiệp nhỏ, dự án cá nhân và môi trường học tập ---
công cụ thường được dùng là các trình cập nhật tự động như Watchtower, vốn chỉ
thực hiện chiến lược Recreate (dừng phiên bản cũ rồi mới chạy phiên bản mới).

% -----------------------------------------------------------------------------
\section{Các tiêu chuẩn và khung tham chiếu}
\label{sec:standards}
\subsection{NIST SSDF (SP 800-218)}
Khung phát triển phần mềm an toàn SSDF phiên bản 1.1~\cite{nist-ssdf} tổ chức
các thực hành thành bốn nhóm: chuẩn bị tổ chức (Prepare the Organization -- PO),
bảo vệ phần mềm (Protect the Software -- PS), tạo ra phần mềm an toàn (Produce
Well-Secured Software -- PW) và phản hồi lỗ hổng (Respond to Vulnerabilities --
RV). Đề tài tập trung vào nhóm PS (bảo vệ tính toàn vẹn của mã nguồn và
artifact, ví dụ ký số và lưu thông tin nguồn gốc) và PW (kiểm tra mã nguồn và
thư viện để phát hiện lỗ hổng một cách tự động).

\subsection{SLSA}
SLSA (Supply-chain Levels for Software Artifacts) phiên bản 1.0~\cite{slsa} định
nghĩa các cấp độ bảo đảm cho quá trình build theo ``Build track'':
\begin{itemize}
  \item \textbf{Build L1:} có thông tin nguồn gốc (provenance) mô tả artifact
  được build như thế nào.
  \item \textbf{Build L2:} build trên nền tảng build được lưu trữ (hosted) và
  provenance được ký bởi nền tảng đó.
  \item \textbf{Build L3:} nền tảng build được gia cố, ngăn các lần build ảnh
  hưởng lẫn nhau và ngăn người dùng giả mạo provenance.
\end{itemize}
Pipeline của đề tài chạy trên runner do GitHub quản lý và ký artifact bằng danh
tính của workflow, hướng tới đáp ứng các yêu cầu của Build L2.

\subsection{OWASP Top 10 CI/CD Security Risks}
Danh sách của OWASP~\cite{owasp-cicd} liệt kê mười nhóm rủi ro đặc thù của hệ
thống CI/CD, từ CICD-SEC-1 đến CICD-SEC-10. Bảng~\ref{tab:owasp-cicd} chọn ra
các rủi ro liên quan trực tiếp đến hệ thống của đề tài.

\begin{table}[!htbp]
  \centering
  \fontsize{11pt}{13.42pt}\selectfont
  \renewcommand{\arraystretch}{1.25}
  \begin{tabularx}{\textwidth}{|>{\raggedright\arraybackslash}p{0.30\textwidth}|>{\raggedright\arraybackslash}X|>{\raggedright\arraybackslash}X|}
    \hline
    \ctuTableHead{Rủi ro} & \ctuTableHead{Biểu hiện trong pipeline niên luận} & \ctuTableHead{Biện pháp trong đề tài} \\
    \hline
    CICD-SEC-1: Insufficient Flow Control Mechanisms & Mã có lỗ hổng vẫn được build và triển khai vì các bước kiểm tra không chặn & Các cổng kiểm soát bắt buộc (Gate 1--3) \\
    \hline
    CICD-SEC-3: Dependency Chain Abuse & Chỉ quét filesystem, kết quả không được dùng để chặn & SCA và quét image với ngưỡng chặn \\
    \hline
    CICD-SEC-4: Poisoned Pipeline Execution & Workflow chạy trên sự kiện \texttt{pull\_request} & Không cho job triển khai trên self-hosted runner chạy với pull request \\
    \hline
    CICD-SEC-6: Insufficient Credential Hygiene & Chưa kiểm tra secret bị commit vào mã nguồn & Quét secret (Gitleaks), dùng token ngắn hạn \\
    \hline
    CICD-SEC-9: Improper Artifact Integrity Validation & Máy chủ chạy bất kỳ image nào có tag \texttt{latest} & Ký image và \texttt{cosign verify} trước khi triển khai \\
    \hline
    CICD-SEC-10: Insufficient Logging and Visibility & Chỉ có log Watchtower & Lưu SBOM, báo cáo quét, log triển khai theo commit SHA \\
    \hline
  \end{tabularx}
  \caption{Các rủi ro OWASP CI/CD liên quan đến đề tài}
  \label{tab:owasp-cicd}
\end{table}
\FloatBarrier

% -----------------------------------------------------------------------------
\section{Nhận xét và khoảng trống nghiên cứu}
\label{sec:gap}
Bảng~\ref{tab:related-compare} so sánh các công trình đã trình bày theo bốn
tiêu chí: trọng tâm, có đề cập bảo mật pipeline hay không, có đề cập chiến lược
triển khai hay không, và phương pháp nghiên cứu.

\begin{table}[!htbp]
  \centering
  \fontsize{11pt}{13.42pt}\selectfont
  \renewcommand{\arraystretch}{1.25}
  \begin{tabularx}{\textwidth}{|>{\raggedright\arraybackslash}p{0.22\textwidth}|>{\raggedright\arraybackslash}X|>{\centering\arraybackslash}p{0.11\textwidth}|>{\centering\arraybackslash}p{0.11\textwidth}|>{\raggedright\arraybackslash}p{0.18\textwidth}|}
    \hline
    \ctuTableHead{Công trình} & \ctuTableHead{Trọng tâm} & \ctuTableHead{Bảo mật pipeline} & \ctuTableHead{Chiến lược triển khai} & \ctuTableHead{Phương pháp} \\
    \hline
    Mohan \& Ben Othmane~\cite{mohan2016} & Hiện trạng nghiên cứu SecDevOps & \checkmark & & Mapping study \\
    \hline
    Myrbakken \& Colomo-Palacios~\cite{myrbakken2017} & Định nghĩa, đặc điểm DevSecOps & \checkmark & & Tổng quan đa nguồn \\
    \hline
    Rajapakse và cs.~\cite{rajapakse2022} & Thách thức, giải pháp DevSecOps & \checkmark & & Tổng quan có hệ thống \\
    \hline
    Enck \& Williams~\cite{enck2022} & Thách thức chuỗi cung ứng & \checkmark & & Thảo luận với tổ chức \\
    \hline
    Ladisa và cs.~\cite{ladisa2023} & Phân loại tấn công chuỗi cung ứng & \checkmark & & Phân loại, khảo sát \\
    \hline
    Newman và cs.~\cite{sigstore-paper} & Ký số keyless (Sigstore) & \checkmark & & Thiết kế hệ thống \\
    \hline
    Shahin và cs.~\cite{shahin2017} & Thực hành CI/CD/CDep & & \checkmark & Tổng quan có hệ thống \\
    \hline
    Humble \& Farley~\cite{humble2010} & Mẫu phân phối liên tục & & \checkmark & Sách chuyên khảo \\
    \hline
    Forsgren và cs.~\cite{forsgren2018} & Chỉ số hiệu quả phân phối & & \checkmark & Khảo sát định lượng \\
    \hline
    \textbf{Tiểu luận này} & \textbf{Pipeline DevSecOps + Blue-Green trên máy chủ đơn} & \checkmark & \checkmark & \textbf{Cài đặt và thực nghiệm so sánh} \\
    \hline
  \end{tabularx}
  \caption{So sánh các công trình liên quan}
  \label{tab:related-compare}
\end{table}
\FloatBarrier

Từ bảng so sánh có thể rút ra các nhận xét sau:
\begin{itemize}
  \item Các công trình về DevSecOps và chuỗi cung ứng chủ yếu là tổng quan tài
  liệu hoặc khảo sát; chúng chỉ ra rằng chi phí thời gian và cảnh báo sai của
  công cụ là rào cản, nhưng ít cung cấp số liệu đo cụ thể cho một pipeline nhỏ.
  \item Các công trình về chiến lược triển khai ít đề cập đến bảo mật chuỗi cung
  ứng, và các hiện thực Blue-Green/Canary phổ biến đều giả định có nền tảng điều
  phối như Kubernetes.
  \item Hai khía cạnh ``bảo mật trước khi triển khai'' và ``an toàn trong khi
  triển khai'' thường được xem xét tách rời, trong khi trên thực tế chúng là hai
  nửa của cùng một quy trình: một image dù đã được kiểm tra kỹ vẫn có thể lỗi
  khi chạy, và một cơ chế triển khai an toàn vẫn vô ích nếu image không đáng tin
  cậy.
\end{itemize}

Từ đó, tiểu luận xác định khoảng trống cần giải quyết là: \textit{thiết kế, cài
đặt và đánh giá định lượng một quy trình kết hợp đồng thời các cổng kiểm soát
bảo mật chuỗi cung ứng và cơ chế triển khai không gián đoạn có tự động quay lui,
trên hạ tầng nhỏ chỉ gồm một máy chủ chạy Docker Compose}. Hai câu hỏi nghiên
cứu RQ1 và RQ2 lần lượt tương ứng với hai khía cạnh này; cơ sở lý thuyết cần
thiết được trình bày ở Chương~\ref{chap:theory}.

\ctutodo{Trước khi nộp: đọc lại bản gốc của từng tài liệu đã trích dẫn để đối
chiếu nội dung tóm tắt (đặc biệt các kết luận của Rajapakse 2022, Enck \&
Williams 2022, Shahin 2017), và bổ sung 1--2 bài báo mới (2023--2026) nếu
GVHD yêu cầu.}

% -----------------------------------------------------------------------------
% CHƯƠNG II
% -----------------------------------------------------------------------------
\chapter{CƠ SỞ LÝ THUYẾT}
\label{chap:theory}
\section{Tổng quan về CI/CD}
CI/CD là một trong những tiến bộ quan trọng trong phát triển phần mềm hiện đại,
giúp thay đổi cách các nhóm DevOps quản lý quá trình phát triển và triển khai
sản phẩm. Ban đầu, các quy trình phát triển phần mềm thường bị chia tách với
quy trình kiểm thử và triển khai, dẫn đến phát hiện lỗi muộn và khó khăn trong
việc duy trì tính ổn định khi tích hợp nhiều thay đổi cùng lúc.

CI/CD (Continuous Integration/Continuous Delivery -- Deployment) là một tập hợp
các phương pháp và công cụ giúp tự động hóa quy trình phát triển, kiểm thử, và
triển khai phần mềm. Mục tiêu chính của CI/CD là đẩy nhanh tốc độ phát hành
phần mềm, giảm lỗi và nâng cao tính ổn định của các phiên bản phát hành
(release).

\ctufigure{image2.png}{1.0}{Quy trình CI/CD}{fig:cicd-process}

\subsection{Tích hợp liên tục}
CI (Continuous Integration) là quá trình liên tục kiểm tra và tích hợp code mới
vào code chính của một dự án. Mục tiêu của CI là đảm bảo rằng bất kỳ thay đổi
nào trong code (chẳng hạn như sửa lỗi hoặc thêm tính năng) đều được kiểm tra
ngay lập tức để phát hiện lỗi trước khi nó gây ảnh hưởng xấu đến hệ thống.

CI hoạt động như một hệ thống kiểm tra tự động. Mỗi khi bạn thay đổi code và
commit lên hệ thống, các bài kiểm tra tự động sẽ được chạy để kiểm tra xem code
có hoạt động đúng và không gây lỗi hay không.

\subsubsection{Những bước cơ bản trong CI}
\begin{enumerate}
  \item \textbf{Functional tests:} Kiểm tra xem phần mềm có hoạt động đúng theo
  mong đợi hay không. Đây thường là các bài kiểm tra tự động mà đội ngũ QA
  (Quality Assurance) phát triển để đảm bảo tính năng hoạt động đúng. VD: Người
  dùng có thể thêm sản phẩm vào giỏ hàng và thực hiện thanh toán thành công hay
  không.
  \item \textbf{Security scans:} Kiểm tra code có lỗ hổng bảo mật nào không,
  chẳng hạn như khả năng bị SQL Injection hay
  \href{https://200lab.io/blog/giai-phap-bao-mat-giao-dien-web-ngan-chan-tan-cong-xss-va-csrf-hieu-qua/}{XSS}
  (Cross-Site Scripting).
  \item \textbf{Code quality scans:} Đảm bảo code tuân thủ các tiêu chuẩn, ví dụ
  như độ dài hàm, cách sử dụng khoảng trắng, và các quy tắc coding style. VD:
  Tiêu chuẩn như đặt tên biến.
  \item \textbf{Performance tests:} Kiểm tra xem code có đáp ứng được yêu cầu về
  hiệu suất không, chẳng hạn như thời gian xử lý yêu cầu hoặc khả năng chịu
  tải.
  \item \textbf{License scanning:} Kiểm tra xem tất cả các thư viện hoặc công cụ
  bạn sử dụng có giấy phép phù hợp hay không, để tránh các vấn đề pháp lý.
  \item \textbf{Fuzz testing:} Gửi các dữ liệu bất thường, như chuỗi ký tự quá
  dài hoặc số không hợp lệ, vào ứng dụng để kiểm tra xem nó có bị crash hay gặp
  lỗi bất ngờ không.
\end{enumerate}

\subsubsection{Lợi ích của việc thiết lập CI}
CI giúp phát hiện sớm các lỗi tiềm ẩn, tiết kiệm rất nhiều thời gian và công
sức, rõ ràng việc sửa một lỗi ngay khi vừa phát hiện sẽ dễ dàng hơn nhiều so
với khi lỗi đó đã tích hợp vào nhiều phần khác của hệ thống.

Ví dụ: Giả sử đội ngũ lập trình đang phát triển một tính năng mới cho ứng dụng.
Thay vì đợi đến cuối quá trình phát triển để kiểm tra, CI cho phép bạn thực
hiện các bài kiểm tra ngay khi lập trình viên viết code. Nhờ đó, nếu có lỗi
xuất hiện, họ có thể phát hiện ngay lập tức và sửa chữa khi vấn đề còn nhỏ,
thay vì đợi đến khi tính năng đã hoàn thành và tích hợp với các phần khác.

CI giúp toàn bộ đội nhóm đều có thể thấy tình trạng hiện tại của phần mềm, tất
cả các thành viên từ lập trình viên, quản lý, đến QA (kiểm thử chất lượng), UX
(trải nghiệm người dùng), và thậm chí cả bộ phận bảo mật có thể theo dõi tiến
trình và trạng thái của dự án. Từ đó, mọi người có thể điều chỉnh công việc của
mình để phù hợp với tình trạng thực tế.

Ví dụ: Khi các bài kiểm tra bảo mật được thực hiện thường xuyên, cả nhóm sẽ
nhận biết được dự án có gặp phải các vấn đề bảo mật mới không. Nếu CI phát hiện
rằng một đoạn code mới làm lộ lỗ hổng bảo mật, thì không chỉ các lập trình viên
mà cả các quản lý dự án cũng biết ngay lập tức. Nhóm sẽ có thể phân bổ nguồn
lực để sửa lỗi hoặc điều chỉnh kế hoạch phát triển.

\subsection{Phân phối/triển khai liên tục}
CD là quá trình đưa code từ nơi bạn viết lên một nơi mà người dùng có thể sử
dụng được, như website hoặc ứng dụng di động. Có 2 loại CD chính:
\begin{itemize}
  \item \textbf{Continuous Delivery (Phân phối liên tục):} Đảm bảo rằng code
  luôn sẵn sàng để được triển khai bất cứ lúc nào. Sau khi code đã vượt qua các
  kiểm tra CI, nó được đặt ở trạng thái sẵn sàng để triển khai. Tuy nhiên, việc
  triển khai có thể được thực hiện thủ công.
  \item \textbf{Continuous Deployment (Triển khai liên tục):} Tự động triển khai
  code ngay khi nó vượt qua tất cả các kiểm tra. Không cần sự can thiệp thủ
  công, hệ thống sẽ tự động đưa code vào môi trường production.
\end{itemize}

\subsubsection{Những bước cơ bản trong CD}
Build: Trước khi triển khai code, thực hiện quá trình build cho các ngôn ngữ
biên dịch (compiled language) như C, C++, Java,~\ldots

Deploy: Triển khai code: đẩy Docker image lên repository, dùng dòng lệnh để đưa
code lên AWS,~\ldots{} Triển khai có thể được tự động hoặc thủ công, tùy thuộc
vào cách bạn thiết lập.

Ví dụ: Khi hoàn thành thêm tính năng mới như giỏ hàng, CD sẽ tự động kiểm tra
và đóng gói tính năng này rồi triển khai lên website thật để khách hàng có thể
sử dụng ngay lập tức.

\subsubsection{Lợi ích của việc thiết lập CD}
Mục tiêu của CD là giúp việc triển khai và phát hành (release) phần mềm trở nên
dễ dàng và có thể diễn ra thường xuyên mà không gặp nhiều rủi ro. Khi CD được
thiết lập, mỗi lần bạn commit code, code có thể tự động được triển khai vào các
môi trường như: review environment, Staging environment hay production
environment.

Các môi trường khác ngoài production (như review và staging) hoạt động giống
như những phép thử, nhằm đảm bảo code chạy ổn định trước khi release chính thức
cho người dùng.

\section{Nền tảng tự động hóa GitHub Actions}
GitHub Actions là nền tảng tự động hóa được tích hợp sẵn trực tiếp trên hệ sinh
thái GitHub. Nền tảng này cho phép tạo các quy trình làm việc (workflow) nhằm
tự động hóa mọi bước trong vòng đời phát triển phần mềm, từ build, test cho đến
deploy ứng dụng ngay khi có sự kiện xảy ra (như push code hoặc tạo pull
request).

Trong kiến trúc của GitHub Actions, một luồng công việc tự động được định nghĩa
thông qua tệp tin cấu hình YAML và bao gồm các thành phần:
\begin{itemize}
  \item \textbf{Events (Sự kiện):} Các trigger kích hoạt quy trình, ví dụ như
  \texttt{push} hoặc \texttt{pull\_request} trên nhánh \texttt{main}.
  \item \textbf{Jobs (Công việc):} Một luồng có thể chạy một hoặc nhiều job (như
  Build, Testing, Security\_Scan, Publish) tuần tự hoặc song song trên các máy
  chủ ảo (runner) do GitHub cung cấp (ví dụ: \texttt{ubuntu-latest}).
  \item \textbf{Steps (Bước):} Từng nhiệm vụ cụ thể bên trong một job, có thể là
  thực thi một lệnh shell hoặc gọi các actions có sẵn (như
  \texttt{actions/checkout} để kéo mã nguồn).
  \item \textbf{Artifacts (Tài nguyên đầu ra):} Chức năng lưu trữ tạm thời các
  gói phần mềm hoặc tệp báo cáo (như file \texttt{.war} hoặc báo cáo bảo mật
  JSON) để chia sẻ giữa các job với nhau.
\end{itemize}

\ctufigure{image3.png}{0.58}{Kiến trúc của GitHub Actions}{fig:gha-architecture}

\section{Công cụ quản lý dự án và kiểm tra mã nguồn Java}
Ứng dụng Web trong đề tài được phát triển dựa trên ngôn ngữ Java, do đó yêu cầu
các công cụ đặc thù để quản lý vòng đời và chất lượng:
\begin{itemize}
  \item Apache Maven: Đây là công cụ quản lý dự án (quản lý thư viện phụ thuộc,
  quản lý vòng đời biên dịch) phổ biến nhất cho Java. Trong hệ thống, Maven
  đóng vai trò tải các thư viện, chạy các kịch bản kiểm thử (\texttt{mvn test})
  và đóng gói mã nguồn thành phần mềm hoàn chỉnh định dạng WAR
  (\texttt{mvn install}).
  \item Checkstyle: Để đảm bảo tính đồng nhất trong phong cách lập trình và phát
  hiện sớm các đoạn code không tuân thủ tiêu chuẩn, công cụ Checkstyle được
  tích hợp vào bước kiểm thử (\texttt{mvn checkstyle:checkstyle}) làm rào chắn
  (coding convention validator).
\end{itemize}

\section{DevSecOps và nguyên tắc Shift-Left Security}
\label{sec:devsecops}
\subsection{Khái niệm DevSecOps}
DevSecOps (Development, Security and Operations) là sự mở rộng của DevOps, trong
đó bảo mật được xem là trách nhiệm chung của cả nhóm phát triển, vận hành và
bảo mật, và được tích hợp vào từng giai đoạn của vòng đời phần mềm thay vì chỉ
là một bước kiểm tra riêng biệt trước khi phát hành~\cite{myrbakken2017}. Về mặt
kỹ thuật, DevSecOps thể hiện qua việc tự động hóa các hoạt động bảo mật (phân
tích mã, kiểm tra thư viện, quét image, ký số, xác minh) và đưa chúng vào chính
pipeline CI/CD, để mọi thay đổi đều được kiểm tra theo cùng một cách, ở cùng một
nơi.

\subsection{Nguyên tắc dịch chuyển sang trái}
Nếu vòng đời phần mềm được biểu diễn từ trái (viết mã) sang phải (vận hành), thì
``dịch chuyển sang trái'' (shift-left) nghĩa là đưa các hoạt động kiểm tra bảo
mật về càng gần thời điểm viết mã càng tốt. Lỗi được phát hiện càng sớm thì
ngữ cảnh càng rõ ràng (lập trình viên vừa viết đoạn mã đó), phạm vi ảnh hưởng
càng nhỏ và chi phí khắc phục càng thấp so với khi lỗi đã được đóng gói, phân
phối và chạy trên môi trường vận hành.

Tuy nhiên, dịch chuyển sang trái không có nghĩa là chỉ kiểm tra ở phía trái.
Như sự cố xz-utils (Mục~\ref{sec:supply-chain}) cho thấy, một số rủi ro chỉ
xuất hiện trong artifact cuối cùng. Do đó đề tài áp dụng mô hình kiểm tra nhiều
lớp: kiểm tra mã nguồn và thư viện ngay sau khi commit, kiểm tra container image
sau khi build, và xác minh chữ ký ngay trước khi triển khai.

\subsection{Cổng kiểm soát bảo mật và chính sách chặn}
Cổng kiểm soát (security gate hoặc quality gate) là một điểm trong pipeline tại
đó kết quả của một hoặc nhiều công cụ kiểm tra được so với một chính sách đã
định trước; nếu vi phạm, pipeline dừng lại và artifact không được chuyển sang
giai đoạn tiếp theo. Một công cụ quét chỉ sinh báo cáo mà không chặn (như cấu
hình Trivy với \texttt{exit-code: 0} trong niên luận) chưa phải là một cổng
kiểm soát.

Chính sách chặn thường dựa trên mức độ nghiêm trọng của lỗ hổng theo thang điểm
CVSS (Common Vulnerability Scoring System)~\cite{cvss} (Bảng~\ref{tab:cvss}).
Một chính sách phổ biến là chặn khi có lỗ hổng mức High hoặc Critical đã có bản
vá.

\begin{table}[!htbp]
  \centering
  \renewcommand{\arraystretch}{1.25}
  \begin{tabular}{|l|c|}
    \hline
    \ctuTableHead{Mức độ} & \ctuTableHead{Điểm CVSS v3} \\
    \hline
    None & 0.0 \\
    \hline
    Low & 0.1 -- 3.9 \\
    \hline
    Medium & 4.0 -- 6.9 \\
    \hline
    High & 7.0 -- 8.9 \\
    \hline
    Critical & 9.0 -- 10.0 \\
    \hline
  \end{tabular}
  \caption{Thang mức độ nghiêm trọng theo CVSS v3}
  \label{tab:cvss}
\end{table}
\FloatBarrier

Khi thiết kế chính sách chặn cần cân nhắc ba vấn đề:
\begin{itemize}
  \item \textbf{Ngưỡng chặn:} ngưỡng quá thấp khiến pipeline liên tục thất bại vì
  các lỗ hổng ít rủi ro, làm nhóm phát triển mất niềm tin vào công cụ; ngưỡng
  quá cao bỏ lọt lỗ hổng nghiêm trọng.
  \item \textbf{Cơ chế ngoại lệ:} một số lỗ hổng chưa có bản vá hoặc không thể
  khai thác trong ngữ cảnh của ứng dụng. Các ngoại lệ cần được ghi rõ mã CVE, lý
  do chấp nhận và thời hạn xem xét lại, đồng thời được lưu trong kho mã để có
  thể truy vết.
  \item \textbf{Chặn khi lỗi (fail-closed):} nếu bản thân công cụ kiểm tra gặp lỗi
  (ví dụ không tải được cơ sở dữ liệu lỗ hổng), pipeline nên dừng thay vì coi
  như kiểm tra đã đạt.
\end{itemize}

% -----------------------------------------------------------------------------
\section{Các kỹ thuật kiểm thử bảo mật tự động}
\label{sec:security-testing}
\subsection{Phân tích mã tĩnh (SAST)}
Kiểm thử bảo mật tĩnh (Static Application Security Testing -- SAST) phân tích mã
nguồn hoặc mã byte mà không cần chạy chương trình, nhằm phát hiện các mẫu lập
trình không an toàn như SQL injection, XSS, đường dẫn tệp không được kiểm soát
hay sử dụng thuật toán mã hóa yếu. Kỹ thuật cốt lõi là phân tích luồng dữ liệu
(taint analysis): theo dõi dữ liệu từ nguồn không tin cậy (source, ví dụ tham số
HTTP) đến các điểm nhạy cảm (sink, ví dụ câu lệnh SQL) mà không đi qua bước làm
sạch (sanitizer).

Ưu điểm của SAST là chạy được rất sớm, chỉ ra chính xác dòng mã có vấn đề. Nhược
điểm là tỉ lệ cảnh báo sai (false positive) có thể cao do công cụ không hiểu đầy
đủ ngữ cảnh chạy, và không phát hiện được lỗi chỉ xuất hiện khi ứng dụng chạy
(lỗi cấu hình máy chủ, lỗi xác thực).

Hai lựa chọn cho ứng dụng Java trong đề tài:
\begin{itemize}
  \item \textbf{CodeQL}~\cite{codeql}: công cụ của GitHub, chuyển mã nguồn thành một
  cơ sở dữ liệu và chạy các truy vấn viết bằng ngôn ngữ QL để tìm lỗ hổng. CodeQL
  tích hợp sẵn với GitHub Actions và hiển thị kết quả trong tab Security của
  kho mã.
  \item \textbf{SpotBugs + FindSecBugs}~\cite{findsecbugs}: SpotBugs phân tích mã byte
  Java để tìm lỗi lập trình; plugin FindSecBugs bổ sung các luật phát hiện lỗ
  hổng bảo mật. Công cụ chạy như một plugin Maven nên dễ tích hợp vào bước build
  hiện có.
\end{itemize}

\subsection{Phân tích thành phần phần mềm (SCA)}
Phân tích thành phần phần mềm (Software Composition Analysis -- SCA) xác định các
thư viện bên thứ ba (kể cả phụ thuộc bắc cầu) mà ứng dụng sử dụng cùng phiên
bản của chúng, sau đó đối chiếu với các cơ sở dữ liệu lỗ hổng công khai như NVD
(National Vulnerability Database) để báo cáo các CVE ảnh hưởng. Với ứng dụng
Java, phần lớn mã thực thi đến từ thư viện, nên SCA thường phát hiện được nhiều
lỗ hổng thực tế hơn SAST.

Đề tài sử dụng \textbf{OWASP Dependency-Check}~\cite{dependency-check}, chạy dưới
dạng plugin Maven. Công cụ suy ra định danh CPE của từng tệp JAR rồi tra cứu
trong NVD; tham số \texttt{failBuildOnCVSS} cho phép làm build thất bại khi có
lỗ hổng vượt ngưỡng điểm CVSS. Do công cụ tải dữ liệu từ API của NVD, cần cấu
hình khóa API và bộ nhớ đệm cơ sở dữ liệu để tránh làm pipeline chậm.

\subsection{Phát hiện secret trong mã nguồn}
Secret là các thông tin bí mật như mật khẩu, khóa API, token truy cập hay khóa
riêng. Khi bị commit vào kho mã, secret vẫn tồn tại trong lịch sử Git ngay cả
khi tệp đã bị xóa ở commit sau. Công cụ \textbf{Gitleaks}~\cite{gitleaks} quét
nội dung và lịch sử commit dựa trên hai kỹ thuật: biểu thức chính quy nhận dạng
định dạng của các loại secret phổ biến (ví dụ token GitHub, khóa AWS), và độ đo
entropy để phát hiện chuỗi ngẫu nhiên có khả năng là khóa. Khi phát hiện secret,
biện pháp đúng là thu hồi và cấp lại secret đó, không chỉ xóa khỏi mã nguồn.

% -----------------------------------------------------------------------------
\section{Quét lỗ hổng container image với Trivy}
\label{sec:trivy}
Trivy~\cite{trivy} là công cụ quét lỗ hổng mã nguồn mở của Aqua Security, hỗ trợ
nhiều loại mục tiêu. Hai chế độ liên quan đến đề tài là:
\begin{itemize}
  \item \textbf{Quét hệ thống tệp (\texttt{fs}):} quét thư mục mã nguồn, nhận dạng
  thư viện qua các tệp khai báo (ví dụ \texttt{pom.xml}). Đây là chế độ được dùng
  trong niên luận.
  \item \textbf{Quét container image (\texttt{image}):} quét image đã build, bao gồm
  các gói của hệ điều hành nền (ví dụ các gói Ubuntu trong image Tomcat), môi
  trường chạy Java và các thư viện nằm trong tệp WAR. Chế độ này phản ánh đúng
  những gì sẽ thực sự chạy trên máy chủ.
\end{itemize}

Trong đề tài, Trivy được chuyển sang quét image theo digest ngay sau bước build,
với tham số \texttt{severity: HIGH,CRITICAL} và \texttt{exit-code: 1} để làm
pipeline thất bại khi có lỗ hổng vượt ngưỡng. Tùy chọn
\texttt{ignore-unfixed} có thể được dùng để bỏ qua các lỗ hổng chưa có bản vá.
Các ngoại lệ được khai báo trong tệp \texttt{.trivyignore} kèm chú thích lý do,
đúng với cơ chế ngoại lệ đã nêu ở Mục~\ref{sec:devsecops}.

% -----------------------------------------------------------------------------
\section{SBOM, ký số artifact và khung SLSA}
\label{sec:sbom-signing}
\subsection{Danh mục thành phần phần mềm (SBOM)}
SBOM (Software Bill of Materials) là danh sách có cấu trúc, đọc được bằng máy,
liệt kê toàn bộ thành phần cấu thành một phần mềm cùng phiên bản, nhà cung cấp,
giấy phép và quan hệ phụ thuộc. Hai định dạng phổ biến là \textbf{SPDX} do Linux
Foundation phát triển (đã được chuẩn hóa thành ISO/IEC 5962:2021)~\cite{spdx} và
\textbf{CycloneDX} do OWASP phát triển~\cite{cyclonedx}. Khi một lỗ hổng mới như
Log4Shell được công bố, tổ chức có SBOM của từng image có thể tra cứu ngay image
nào bị ảnh hưởng mà không cần quét lại. Đề tài dùng công cụ Syft để sinh SBOM
cho mỗi image.

\subsection{Ký số container image với Sigstore/Cosign}
Container image trong registry được định danh bởi digest (giá trị băm SHA-256
của manifest), nên nội dung của một digest không thể bị thay đổi. Tuy nhiên tag
(như \texttt{latest}) có thể được trỏ sang digest khác bởi bất kỳ ai có quyền
ghi vào registry. Ký số cho phép máy chủ xác minh rằng image cần chạy được tạo
ra bởi đúng pipeline của dự án.

Cosign, thuộc dự án Sigstore~\cite{sigstore-paper, sigstore}, hỗ trợ ký theo cơ chế
``keyless'' (Hình~\ref{fig:cosign}):
\begin{enumerate}
  \item Workflow GitHub Actions yêu cầu một token OIDC từ GitHub; token này chứa
  danh tính của workflow (kho mã, tệp workflow, nhánh).
  \item Cosign gửi token đến \textbf{Fulcio} (cơ quan cấp chứng chỉ của Sigstore)
  để nhận một chứng chỉ ngắn hạn gắn với danh tính đó và một cặp khóa tạm thời.
  \item Cosign ký digest của image, lưu chữ ký vào registry bên cạnh image, và
  ghi bản ghi vào \textbf{Rekor} -- nhật ký minh bạch công khai, không thể sửa.
  \item Khi xác minh, \texttt{cosign verify} kiểm tra chữ ký, chứng chỉ và bản
  ghi Rekor, đồng thời kiểm tra danh tính người ký có khớp với workflow mong đợi
  (tham số \texttt{-{}-certificate-identity} và
  \texttt{-{}-certificate-oidc-issuer}).
\end{enumerate}
Vì khóa ký chỉ tồn tại trong thời gian chạy job, không có khóa bí mật dài hạn
nào cần lưu trữ hay có thể bị lộ.

\begin{figure}[!htbp]
  \centering
  \begin{tikzpicture}[
      box/.style={draw, rounded corners, align=center, minimum height=11mm,
        text width=32mm, font=\small},
      lbl/.style={font=\footnotesize, align=center},
      >=stealth]
    \node[box] (gha) at (0,0) {Workflow GitHub Actions};
    \node[box] (fulcio) at (8.5,0) {Fulcio (cơ quan cấp chứng chỉ)};
    \node[box] (ghcr) at (0,-3.4) {GHCR\\ (image + chữ ký)};
    \node[box] (rekor) at (8.5,-3.4) {Rekor (nhật ký minh bạch)};
    \node[box] (srv) at (0,-6.6) {Máy chủ triển khai};
    \draw[->, thick] ([yshift=2mm]gha.east) -- node[lbl, above] {(1) token OIDC} ([yshift=2mm]fulcio.west);
    \draw[->, thick] ([yshift=-2mm]fulcio.west) -- node[lbl, below] {(2) chứng chỉ ngắn hạn} ([yshift=-2mm]gha.east);
    \draw[->, thick] (gha) -- node[lbl, left] {(3) đẩy chữ ký} (ghcr);
    \draw[->, thick] (gha.south east) -- node[lbl, below, sloped] {(3) ghi bản ghi chữ ký} (rekor.north west);
    \draw[->, thick, dashed] (srv) -- node[lbl, left] {(4) cosign verify} (ghcr);
    \draw[->, thick, dashed] (srv.east) -| node[lbl, below, pos=0.25] {(4) đối chiếu bản ghi} (rekor.south);
  \end{tikzpicture}
  \caption{Luồng ký và xác minh image theo cơ chế keyless của Sigstore}
  \label{fig:cosign}
\end{figure}
\FloatBarrier

\subsection{Attestation và thông tin nguồn gốc}
Ngoài chữ ký trên image, Cosign còn có thể tạo \textit{attestation}: một khẳng
định đã được ký theo định dạng in-toto, gắn một tài liệu (ví dụ SBOM hoặc thông
tin nguồn gốc build -- provenance) với digest của image. Nhờ đó, máy chủ có thể
xác minh không chỉ ``ai đã ký image'' mà còn ``image này chứa những thành phần
nào'' và ``được build từ commit nào''. Như đã trình bày ở
Mục~\ref{sec:standards}, việc build trên runner do GitHub quản lý kết hợp với
provenance được ký là các yêu cầu của SLSA Build L2~\cite{slsa}.

% -----------------------------------------------------------------------------
\section{Công nghệ ảo hóa (Docker)}
Docker là nền tảng phần mềm cho phép các nhà phát triển xây dựng, kiểm thử và
triển khai các ứng dụng thông qua công nghệ ảo hóa ở dạng container. Docker
đóng gói ứng dụng Web Java cùng toàn bộ môi trường chạy thực tế (runtime, thư
viện hệ thống) thành các đơn vị chuẩn hóa gọi là Container. Nhờ vậy, ứng dụng
có thể chạy đồng nhất trên bất kỳ hạ tầng nào mà không gặp tình trạng xung đột
môi trường. Trong dự án, Dockerfile của ứng dụng dùng image nền
\texttt{tomcat:10-jdk21} và chỉ sao chép tệp WAR đã được build ở bước CI
(nguyên tắc ``build once''), nên image không cần chứa Maven hay mã nguồn.

Docker Compose cho phép mô tả nhiều container liên quan (ứng dụng, cơ sở dữ
liệu, bộ nhớ đệm, hàng đợi thông điệp, proxy) trong một tệp YAML và quản lý
chúng như một đơn vị. Các container trong cùng một dự án Compose được đặt trong
một mạng nội bộ và gọi nhau bằng tên dịch vụ; đây là cơ sở để Nginx chuyển
hướng lưu lượng giữa hai dịch vụ Blue và Green.

\section{GitHub Container Registry (GHCR)}
GitHub Container Registry và Docker Hub là hai nền tảng lưu trữ, quản lý và
phân phối container image phổ biến. Tuy nhiên, thay vì sử dụng AWS ECR hay
Docker Hub, đề tài lựa chọn GHCR vì đây là dịch vụ được tích hợp sẵn trong hệ
sinh thái GitHub.

Việc sử dụng GHCR mang lại những lợi ích cốt lõi cho hệ thống:
\begin{itemize}
  \item Cho phép nhà phát triển lưu trữ Docker image ngay tại nơi họ đang quản lý
  mã nguồn.
  \item Kết nối cực kỳ liền mạch và có độ tin cậy cao với GitHub Actions.
  \item Đảm bảo an toàn bảo mật thông qua cơ chế phân quyền chi tiết (granular
  permissions) và việc xác thực tự động không cần chia sẻ mật khẩu tĩnh
  (secretless authentication) thông qua token mặc định của GitHub.
\end{itemize}

\section{Watchtower và cơ chế tự động cập nhật container}
Watchtower là công cụ hỗ trợ tự động cập nhật Docker container. Công cụ này
chạy như một container trên máy chủ triển khai và định kỳ kiểm tra Docker
registry để phát hiện image mới của các container được theo dõi.

Trong hệ thống của niên luận, Watchtower được cấu hình trên máy ảo Ubuntu có cài
Docker Engine. Sau khi GitHub Actions xuất bản Docker image mới lên GHCR,
Watchtower sẽ xác thực với GHCR, tải image mới về máy chủ, dừng container ứng
dụng cũ và khởi động lại container mới từ image vừa tải. Luồng triển khai tổng
quát như sau:

\begin{center}
  \begin{minipage}{0.92\textwidth}
    \raggedright
    Push source code $\rightarrow$ GitHub Actions
    build/test/checkstyle/security scan $\rightarrow$ Docker Build
    $\rightarrow$ Push Docker image lên GHCR $\rightarrow$ Watchtower phát hiện
    image mới $\rightarrow$ Pull image $\rightarrow$ Dừng container cũ
    $\rightarrow$ Khởi động container mới.
  \end{minipage}
\end{center}

Cơ chế này đơn giản và phù hợp với máy chủ đơn, nhưng có ba hạn chế quan trọng
đối với mục tiêu của tiểu luận:
\begin{itemize}
  \item Container cũ bị dừng trước khi container mới sẵn sàng, nên ứng dụng bị
  gián đoạn trong suốt thời gian Tomcat khởi động (chiến lược Recreate, xem
  Mục~\ref{sec:deploy-strategies}).
  \item Watchtower không chờ ứng dụng trong container mới thực sự phục vụ được
  yêu cầu và không tự động quay lại phiên bản cũ khi phiên bản mới bị lỗi.
  \item Watchtower theo dõi một tag (ví dụ \texttt{latest}) và không xác minh chữ
  ký, nên sẽ chạy bất kỳ image nào được gắn tag đó trong registry.
\end{itemize}
Trong tiểu luận, Watchtower được giữ lại làm phương án đối chứng (baseline) cho
các thực nghiệm ở Chương~\ref{chap:evaluation}.

% -----------------------------------------------------------------------------
\section{Các chiến lược triển khai ứng dụng}
\label{sec:deploy-strategies}
Chiến lược triển khai quyết định cách phiên bản mới thay thế phiên bản đang chạy
và ảnh hưởng trực tiếp đến thời gian gián đoạn, tài nguyên cần thiết và khả năng
quay lui~\cite{humble2010}.

\textbf{Recreate.} Dừng toàn bộ phiên bản cũ, sau đó khởi động phiên bản mới.
Đơn giản và không bao giờ có hai phiên bản chạy song song, nhưng luôn có một
khoảng gián đoạn bằng thời gian khởi động của phiên bản mới. Muốn quay lui phải
triển khai lại phiên bản cũ, cũng chịu thêm một khoảng gián đoạn. Watchtower
thực hiện chiến lược này.

\textbf{Rolling Update.} Khi ứng dụng chạy nhiều bản sao (replica) sau bộ cân
bằng tải, các bản sao được thay thế lần lượt, nên luôn còn bản sao phục vụ yêu
cầu. Chiến lược này là mặc định của Kubernetes, nhưng đòi hỏi nhiều bản sao và
trong lúc cập nhật có cả hai phiên bản cùng phục vụ người dùng.

\textbf{Blue-Green.} Duy trì hai môi trường giống nhau, ``Blue'' đang phục vụ và
``Green'' rảnh. Phiên bản mới được triển khai lên môi trường rảnh, kiểm tra cho
đến khi sẵn sàng, sau đó bộ định tuyến chuyển toàn bộ lưu lượng sang. Môi trường
cũ được giữ nguyên để có thể chuyển lưu lượng trở lại gần như tức
thời~\cite{fowler-bluegreen} (Hình~\ref{fig:bluegreen}).

\textbf{Canary.} Phiên bản mới nhận một phần nhỏ lưu lượng (ví dụ 5\%), các chỉ
số lỗi và độ trễ được theo dõi; nếu ổn định thì tăng dần tỉ lệ. Canary giảm
phạm vi ảnh hưởng của lỗi nhưng cần hệ thống giám sát và cơ chế phân chia lưu
lượng theo trọng số, và thời gian triển khai kéo dài hơn.

\begin{figure}[!htbp]
  \centering
  \begin{tikzpicture}[
      box/.style={draw, rounded corners, align=center, minimum height=10mm,
        text width=36mm, font=\small},
      >=stealth]
    \node[box] (user) at (0,0) {Người dùng};
    \node[box] (nginx) at (0,-2) {Nginx (reverse proxy)};
    \node[box, fill=blue!10] (blue) at (-3.3,-4.6) {vproapp-blue\\ phiên bản N (đang phục vụ)};
    \node[box, fill=green!10] (green) at (3.3,-4.6) {vproapp-green\\ phiên bản N+1 (đang kiểm tra)};
    \node[box, text width=70mm] (shared) at (0,-7.2) {MySQL, Memcached, RabbitMQ (dùng chung)};
    \draw[->, thick] (user) -- (nginx);
    \draw[->, very thick] (nginx) -- node[left, font=\footnotesize, pos=0.4] {100\% lưu lượng} (blue);
    \draw[->, thick, dashed] (nginx) -- node[right, font=\footnotesize, pos=0.4] {sau khi sẵn sàng} (green);
    \draw[->] (blue) -- (shared);
    \draw[->] (green) -- (shared);
  \end{tikzpicture}
  \caption{Mô hình triển khai Blue-Green với Nginx}
  \label{fig:bluegreen}
\end{figure}
\FloatBarrier

Bảng~\ref{tab:deploy-compare} so sánh bốn chiến lược trong bối cảnh của đề tài.

\begin{table}[!htbp]
  \centering
  \fontsize{11pt}{13.42pt}\selectfont
  \renewcommand{\arraystretch}{1.25}
  \begin{tabularx}{\textwidth}{|>{\raggedright\arraybackslash}p{0.17\textwidth}|>{\raggedright\arraybackslash}X|>{\raggedright\arraybackslash}X|>{\raggedright\arraybackslash}X|>{\raggedright\arraybackslash}X|}
    \hline
    \ctuTableHead{Tiêu chí} & \ctuTableHead{Recreate} & \ctuTableHead{Rolling} & \ctuTableHead{Blue-Green} & \ctuTableHead{Canary} \\
    \hline
    Gián đoạn & Có & Không & Không & Không \\
    \hline
    Tài nguyên thêm & Không & Thêm bản sao tạm thời & Gấp đôi tầng ứng dụng khi triển khai & Thêm bản sao canary \\
    \hline
    Quay lui & Triển khai lại, có gián đoạn & Cập nhật ngược lần lượt & Chuyển lưu lượng, gần như tức thời & Rút lưu lượng khỏi canary \\
    \hline
    Hai phiên bản cùng phục vụ & Không & Có & Không & Có \\
    \hline
    Độ phức tạp trên máy chủ đơn & Thấp & Cần nhiều bản sao và cân bằng tải & Trung bình & Cao, cần giám sát theo chỉ số \\
    \hline
  \end{tabularx}
  \caption{So sánh các chiến lược triển khai}
  \label{tab:deploy-compare}
\end{table}
\FloatBarrier

Với hạ tầng một máy chủ chạy Docker Compose, Blue-Green là lựa chọn cân bằng: loại
bỏ gián đoạn và có khả năng quay lui tức thời, trong khi chỉ cần thêm một
container ứng dụng trong lúc triển khai và một bộ định tuyến đã có sẵn (Nginx).
Việc chuyển lưu lượng được thực hiện bằng cách thay đổi cấu hình upstream rồi
chạy \texttt{nginx -s reload}: tiến trình master của Nginx khởi động các worker
mới với cấu hình mới, còn các worker cũ ngừng nhận kết nối mới và kết thúc sau
khi xử lý xong các yêu cầu đang dở~\cite{nginx-control}, nên không có yêu cầu nào
bị cắt ngang.

Hạn chế chung của Blue-Green là cơ sở dữ liệu thường được dùng chung giữa hai
môi trường; nếu phiên bản mới thay đổi lược đồ dữ liệu theo cách không tương
thích ngược thì việc quay lui không còn an toàn. Vấn đề này nằm ngoài phạm vi
của đề tài, như đã nêu trong phần Giới thiệu.

% -----------------------------------------------------------------------------
\section{Health check và cơ chế tự động quay lui}
\label{sec:healthcheck}
\subsection{Kiểm tra tình trạng ứng dụng}
Một container ở trạng thái ``đang chạy'' chưa chắc ứng dụng bên trong đã phục vụ
được yêu cầu: Tomcat có thể còn đang triển khai tệp WAR, hoặc ứng dụng khởi động
nhưng không kết nối được cơ sở dữ liệu. Kiểm tra tình trạng (health check) giải
quyết vấn đề này bằng cách định kỳ gửi một yêu cầu thử đến ứng dụng. Có thể phân
biệt hai loại kiểm tra:
\begin{itemize}
  \item \textbf{Liveness:} tiến trình còn sống và phản hồi hay không; dùng để quyết
  định có khởi động lại container hay không.
  \item \textbf{Readiness:} ứng dụng đã sẵn sàng nhận lưu lượng hay chưa; dùng để
  quyết định có chuyển lưu lượng đến hay không. Với Blue-Green, đây là loại kiểm
  tra quyết định thời điểm chuyển đổi.
\end{itemize}

Docker hỗ trợ khai báo health check qua chỉ thị \texttt{HEALTHCHECK} trong
Dockerfile hoặc khóa \texttt{healthcheck} trong Compose, với các tham số
\texttt{test} (lệnh kiểm tra), \texttt{interval}, \texttt{timeout},
\texttt{retries} và \texttt{start\_period} (khoảng thời gian khởi động mà lỗi chưa
bị tính)~\cite{docker-healthcheck}. Container sẽ có trạng thái \texttt{starting},
\texttt{healthy} hoặc \texttt{unhealthy}. Cần lưu ý rằng Docker Engine (không dùng
Swarm) chỉ đánh dấu trạng thái chứ không tự khởi động lại hay thay thế container
\texttt{unhealthy}; hành động phải do kịch bản triển khai đảm nhiệm.

\subsection{Tiêu chí xác định phiên bản lỗi}
Đề tài sử dụng hai tiêu chí, tương ứng với hai thời điểm:
\begin{enumerate}
  \item \textbf{Trước khi chuyển lưu lượng:} container mới phải đạt trạng thái
  \texttt{healthy} (trang đăng nhập trả về mã HTTP 200) trong một khoảng thời gian
  tối đa; nếu không, việc triển khai bị hủy và lưu lượng vẫn ở môi trường cũ.
  \item \textbf{Sau khi chuyển lưu lượng:} trong một cửa sổ quan sát (ví dụ 60
  giây), kịch bản tiếp tục kiểm tra phiên bản mới; nếu số lần thất bại vượt
  ngưỡng, lưu lượng được chuyển trở lại môi trường cũ.
\end{enumerate}

\subsection{Quy trình quay lui}
Nhờ môi trường cũ vẫn chạy trong suốt cửa sổ quan sát, quay lui chỉ gồm thao tác
ghi lại cấu hình upstream trỏ về môi trường cũ và nạp lại Nginx, không cần tải
lại image hay khởi động lại Tomcat. Chỉ sau khi cửa sổ quan sát kết thúc thành
công, môi trường cũ mới được dừng để giải phóng tài nguyên. Ngoài ra, việc gắn
tag image theo commit SHA (thay vì \texttt{latest}) bảo đảm luôn xác định được
chính xác phiên bản nào đang chạy và phiên bản nào là bản ổn định trước đó.

% -----------------------------------------------------------------------------
\section{Các chỉ số đánh giá}
\label{sec:metrics}
\subsection{Các chỉ số DORA}
Chương trình nghiên cứu DORA~\cite{dora, forsgren2018} đề xuất bốn chỉ số then
chốt, chia thành hai nhóm:
\begin{itemize}
  \item \textbf{Thông lượng (throughput):} tần suất triển khai (deployment
  frequency) và thời gian từ khi commit đến khi chạy trên môi trường vận hành
  (lead time for changes).
  \item \textbf{Độ ổn định (stability):} tỉ lệ thay đổi gây lỗi (change failure
  rate) và thời gian khôi phục dịch vụ sau một lần triển khai lỗi (time to
  restore service; các báo cáo gần đây gọi là failed deployment recovery time).
\end{itemize}
Các chỉ số DORA được thiết kế để đo một tổ chức trong thời gian dài, nên không
thể đo đầy đủ trong một môi trường thực nghiệm. Đề tài dùng chúng làm định
hướng: các cổng bảo mật ảnh hưởng đến \textit{lead time}, còn Blue-Green và quay
lui tự động ảnh hưởng đến \textit{time to restore}.

\subsection{Các chỉ số dùng trong thực nghiệm}
Bảng~\ref{tab:metrics} định nghĩa các chỉ số được đo trong
Chương~\ref{chap:evaluation}.

\begin{table}[!htbp]
  \centering
  \fontsize{11pt}{13.42pt}\selectfont
  \renewcommand{\arraystretch}{1.25}
  \begin{tabularx}{\textwidth}{|>{\raggedright\arraybackslash}p{0.25\textwidth}|>{\raggedright\arraybackslash}X|>{\centering\arraybackslash}p{0.10\textwidth}|}
    \hline
    \ctuTableHead{Chỉ số} & \ctuTableHead{Định nghĩa} & \ctuTableHead{Thực nghiệm} \\
    \hline
    Tỉ lệ phát hiện & Số kịch bản lỗi cố ý bị chặn trên tổng số kịch bản & E1 \\
    \hline
    Cảnh báo sai & Số cảnh báo không tương ứng với lỗ hổng thực sự & E1 \\
    \hline
    Thời gian pipeline & Thời gian từ khi workflow được kích hoạt đến khi image được xuất bản, và thời gian của từng job & E2 \\
    \hline
    Số yêu cầu lỗi & Số yêu cầu HTTP thất bại hoặc trả về mã 5xx trong lúc triển khai & E3 \\
    \hline
    Thời gian gián đoạn & Khoảng thời gian liên tục dài nhất mà mọi yêu cầu đều thất bại & E3 \\
    \hline
    Độ trễ p95 & Giá trị mà 95\% yêu cầu có thời gian phản hồi nhỏ hơn hoặc bằng & E3 \\
    \hline
    MTTR & Thời gian từ khi phiên bản lỗi bắt đầu nhận lưu lượng (hoặc bắt đầu triển khai) đến khi dịch vụ phục vụ bình thường trở lại & E4 \\
    \hline
  \end{tabularx}
  \caption{Các chỉ số đánh giá trong thực nghiệm}
  \label{tab:metrics}
\end{table}
\FloatBarrier

Do thời gian chạy trên runner của GitHub và thời gian khởi động container dao
động giữa các lần chạy, mỗi thực nghiệm định lượng được lặp lại nhiều lần; kết
quả được báo cáo dưới dạng giá trị trung bình kèm độ lệch chuẩn, và trung vị khi
phân bố bị lệch. Công cụ k6~\cite{k6} được dùng để tạo tải HTTP liên tục và ghi
nhận các chỉ số E3.

% -----------------------------------------------------------------------------
% CHƯƠNG III
% -----------------------------------------------------------------------------
\chapter{PHÂN TÍCH VÀ THIẾT KẾ HỆ THỐNG}
\label{chap:design}
\section{Hiện trạng hệ thống từ niên luận}
\subsection{Phân tích kiến trúc hệ thống ứng dụng Web mẫu}
Ứng dụng web mẫu trong đề tài (VProfile) được viết bằng ngôn ngữ lập trình Java
và được thiết kế theo mô hình đa tầng (Multi-tier) bao gồm các thành phần:
\begin{itemize}
  \item Tầng ứng dụng web (Web Tier): Sử dụng Nginx làm reverse proxy và cân bằng
  tải (load balancer).
  \item Tầng Ứng dụng (Application Tier): Được viết bằng ngôn ngữ Java và
  Framework Spring, chạy trên máy chủ Apache Tomcat (port 8080). Đồng thời đóng
  gói phần mềm đầu ra của tầng này là một tệp định dạng \texttt{.war}.
  \item Tầng Cơ sở dữ liệu và Dịch vụ hỗ trợ: Bao gồm MySQL 8.0, Memcached
  (caching), và RabbitMQ (Message Broker).
  \item Hệ thống CI/CD trong đề tài sẽ tập trung xử lý tự động hóa cho Tầng Ứng
  dụng web (Java Spring MVC), đảm bảo mã nguồn được kiểm thử, đóng gói thành
  file \texttt{.war}, và cuối cùng xây dựng thành Docker image để sẵn sàng tích
  hợp với các tầng dịch vụ khác.
\end{itemize}

\ctufigure{image4.png}{1.0}{Kiến trúc Web App Java}{fig:webapp-architecture}
\subsection{Thiết kế luồng công việc tự động hóa (Pipeline Design)}
Để tối ưu hóa thời gian chạy và đảm bảo tính chặt chẽ, pipeline được chia thành
5 công việc (Jobs) chính, hoạt động kết hợp giữa tuần tự và song song:
\begin{enumerate}
  \item Job CI (Build, Test \& Package): Chạy đầu tiên để kiểm thử mã nguồn và
  đóng gói thành một artifact (\texttt{.war}).
  \item Job Security\_Scan: Hoạt động song song với Job CI nhằm đánh giá bảo mật
  mà không làm kéo dài thời gian chờ đợi.
  \item Job BUILD\_AMD64 \& BUILD\_ARM64: Chỉ được kích hoạt khi cả hai khâu
  kiểm thử và bảo mật đều thành công. Hai job này chạy song song trên các hạ
  tầng khác nhau để đóng gói Docker cho cả hai kiến trúc vi xử lý (x86\_64 và
  ARM64).
  \item Job PUBLISH\_TO\_GHCR: Hoạt động cuối cùng, gộp kết quả của hai tiến
  trình build trên thành một container đa kiến trúc (Multi-arch) và phân phối
  lên kho lưu trữ của GitHub (GHCR).
  \item Sau khi image được xuất bản thành công, dịch vụ Watchtower trên máy ảo
  Ubuntu sẽ tự động kiểm tra, tải phiên bản image mới và khởi động lại container
  ứng dụng. Qua đó, hệ thống thực hiện Continuous Deployment trong phạm vi môi
  trường thực nghiệm.
\end{enumerate}

\section{Phân tích hạn chế và rủi ro}
\subsection{Hạn chế của pipeline hiện tại}
\ctutodo{Liệt kê có dẫn chứng: Trivy \texttt{exit-code: 0}; Checkstyle không
chặn; \texttt{Tests run: 0}; không có SAST, SBOM, chữ ký; Watchtower dùng tag
\texttt{latest}, dừng rồi mới chạy container mới, không health check, không
rollback.}

\subsection{Mô hình hóa mối đe dọa theo STRIDE}
\ctutodo{Vẽ sơ đồ luồng dữ liệu (lập trình viên $\rightarrow$ GitHub $\rightarrow$
runner $\rightarrow$ GHCR $\rightarrow$ máy chủ). Lập bảng: thành phần, mối đe
dọa STRIDE, ví dụ, biện pháp giảm thiểu trong thiết kế mới.}

\section{Yêu cầu của hệ thống mới}
\subsection{Yêu cầu chức năng}
\ctutodo{Ví dụ: FR1 pipeline thất bại khi có lỗ hổng HIGH/CRITICAL; FR2 mỗi image
có SBOM và chữ ký; FR3 máy chủ từ chối image chưa ký; FR4 triển khai không gián
đoạn; FR5 tự động rollback khi health check thất bại.}
\subsection{Yêu cầu phi chức năng}
\ctutodo{Ví dụ: thời gian pipeline tăng không quá X\%; thời gian rollback dưới Y
giây; mọi phiên bản truy vết được qua commit SHA.}

\section{Thiết kế pipeline DevSecOps}
\ctutodo{Sơ đồ pipeline mới (TikZ hoặc hình): Build \& Unit test (JaCoCo)
$\parallel$ SAST $\parallel$ SCA $\parallel$ Secret scan $\rightarrow$ Gate 1
$\rightarrow$ Build image amd64/arm64 $\rightarrow$ Trivy image scan
$\rightarrow$ Gate 2 $\rightarrow$ Push theo digest, SBOM, ký Cosign
$\rightarrow$ Gate 3 $\rightarrow$ Deploy. Bảng: gate, công cụ, điều kiện chặn,
ngoại lệ.}

\section{Thiết kế cơ chế triển khai Blue-Green}
\subsection{Kiến trúc triển khai}
\ctutodo{Sơ đồ: Nginx $\rightarrow$ vproapp-blue / vproapp-green; dịch vụ dùng
chung (MySQL, Memcached, RabbitMQ). Thay \texttt{container\_name} cố định bằng
hai service; upstream Nginx sinh từ template.}
\subsection{Luồng triển khai và quay lui}
\ctutodo{Sơ đồ trạng thái / lưu đồ của \texttt{deploy.sh}: xác minh chữ ký
$\rightarrow$ khởi động màu rảnh $\rightarrow$ chờ health check $\rightarrow$
chuyển upstream (\texttt{nginx -s reload}) $\rightarrow$ theo dõi N giây
$\rightarrow$ giữ hoặc rollback.}
\subsection{Cơ chế kích hoạt triển khai}
\ctutodo{Lý do chọn self-hosted runner trên máy ảo (không mở SSH ra Internet, mô
hình pull giống Watchtower nên so sánh công bằng); phân quyền của runner.}

% -----------------------------------------------------------------------------
% CHƯƠNG IV
% -----------------------------------------------------------------------------
\chapter{CÀI ĐẶT HỆ THỐNG}
\label{chap:implementation}
\section{Cài đặt quy trình CI (Continuous Integration)}
Quá trình cài đặt giải pháp được thực hiện thông qua tệp tin cấu hình YAML
(\texttt{.github/workflows/main.yml}) đặt tại kho lưu trữ mã nguồn.

\subsection{Tạo access token trên GitHub}
Tạo \textbf{Personal access tokens} để đăng nhập vào GHCR (GitHub Container
Registry).

Các quyền cần thiết:
\begin{itemize}
  \item Quyền quản lý repo
  \item Quyền cập nhật workflow
  \item Quyền ghi và đọc packages
\end{itemize}

\ctufigure{image5.png}{1.0}{Tạo personal access token (classic)}{fig:pat}

\subsection{Cấu hình sự kiện kích hoạt (Triggers)}
Hệ thống được thiết lập cơ chế tự động kích hoạt (trigger) khi có sự kiện đẩy
mã nguồn (\texttt{push}) hoặc tạo yêu cầu gộp mã (\texttt{pull\_request}) vào
nhánh chính (\texttt{main}). Nền tảng được cấp quyền \texttt{contents: read} để
sao chép mã nguồn và \texttt{packages: write} để đẩy các container lên GitHub
Container Registry (GHCR).

\ctufigure{image6.png}{0.42}{Cấu hình sự kiện kích hoạt}{fig:triggers}

\subsection{Cài đặt giai đoạn Tích hợp liên tục (CI)}
Giai đoạn Tích hợp liên tục chịu trách nhiệm biên dịch và kiểm thử ứng dụng
Java.
\begin{itemize}
  \item Thiết lập môi trường: Khởi tạo máy chủ ảo \texttt{ubuntu-latest}, cài
  đặt JDK 17 (Temurin). Đặc biệt, cơ chế bộ nhớ đệm (\texttt{cache: 'maven'})
  được kích hoạt để lưu lại các thư viện phụ thuộc, giúp giảm thiểu thời gian
  tải lại các package trong các lần chạy sau.
  \item Biên dịch và Kiểm thử: Lệnh \texttt{mvn install} được thực thi để
  compile, test, và package ra tệp \texttt{.war}.
  \item Kiểm tra mã nguồn: Lệnh \texttt{mvn checkstyle:checkstyle} được gọi để
  đảm bảo mã nguồn tuân thủ các quy tắc lập trình (coding convention).
  \item Lưu trữ Artifact đầu ra: Gói phần mềm \texttt{vprofile-v2.war} được
  trích xuất và tải lên bộ nhớ tạm của GitHub Actions (thời hạn 7 ngày) để tiết
  kiệm tài nguyên.
\end{itemize}

\ctufigure{image7.png}{1.0}{Job CI}{fig:job-ci}

\subsection{Cài đặt giai đoạn Đánh giá bảo mật DevSecOps}
Tuân thủ nguyên tắc ``Shift-Left Security'', công việc Security\_Scan được cấu
hình chạy song song độc lập với tiến trình CI.

Hệ thống sử dụng action \texttt{aquasecurity/trivy-action} để phân tích tĩnh
toàn bộ hệ thống tập tin (filesystem), tìm kiếm các lỗ hổng của thư viện và hệ
điều hành dựa trên dữ liệu CVE. Kết quả báo cáo được xuất ra dưới định dạng
JSON làm cơ sở đánh giá chất lượng.

\ctufigure{image8.png}{1.0}{Job Security\_Scan}{fig:job-security}

\subsection{Cài đặt giai đoạn Đóng gói đa kiến trúc (Multi-Arch Docker Build)}
Thay vì tự biên dịch lại mã nguồn bên trong Docker (gây dư thừa tài nguyên), hệ
thống áp dụng nguyên tắc ``Build Once'' (Chỉ biên dịch một lần).

Các Job đóng gói sẽ tải tệp \texttt{.war} từ bước CI về. Dockerfile chỉ cần
thực hiện lệnh \texttt{COPY} tệp tin này vào máy chủ Tomcat.

Để hỗ trợ tính tương thích cho nhiều nền tảng, mã nguồn được build đồng thời
cho hai kiến trúc:
\begin{itemize}
  \item AMD64: Thực thi trên runner mặc định \texttt{ubuntu-latest}.
\end{itemize}

\ctufigure{image9.png}{1.0}{Job BUILD\_AMD64}{fig:job-amd64}

\begin{itemize}
  \item ARM64: Thực thi trực tiếp trên hạ tầng Native ARM
  (\nolinkurl{ubuntu-24.04-arm}). Việc không sử dụng giả lập QEMU (Quick EMUlator)
  giúp tốc độ build phần mềm tăng lên đáng kể.
\end{itemize}

\ctufigure{image10.png}{1.0}{Job BUILD\_ARM64}{fig:job-arm64}

\begin{itemize}
  \item Thay vì gán nhãn \texttt{latest} ngay lúc này, cả hai Job đều đẩy (push)
  image lên GHCR dưới dạng chuỗi mã băm (Digest) để tránh xung đột dữ liệu.
\end{itemize}

\section{Cài đặt các cổng kiểm soát bảo mật}
\subsection{Bổ sung kiểm thử tự động và đo độ phủ}
\ctutodo{Viết unit test cho controller/service, cấu hình JaCoCo; chụp kết quả
\texttt{Tests run} khác 0 và báo cáo độ phủ.}
\subsection{Cấu hình Checkstyle chặn build}
\ctutodo{\texttt{failOnViolation}, ngưỡng vi phạm.}
\subsection{Phân tích mã tĩnh}
\ctutodo{Job CodeQL hoặc SpotBugs + FindSecBugs; đoạn YAML; kết quả.}
\subsection{Phân tích thư viện phụ thuộc}
\ctutodo{Job OWASP Dependency-Check, ngưỡng \texttt{failBuildOnCVSS}; báo cáo.}
\subsection{Quét secret}
\ctutodo{Job Gitleaks.}
\subsection{Quét container image}
\ctutodo{Chuyển Trivy sang quét image theo digest, \texttt{exit-code: 1},
\texttt{severity: HIGH,CRITICAL}, \texttt{.trivyignore} kèm lý do.}
\subsection{Sinh SBOM và ký số image}
\ctutodo{Syft sinh SBOM; \texttt{cosign sign} keyless, \texttt{cosign attest}
SBOM; quyền \texttt{id-token: write}. Gắn tag theo commit SHA.}

\section{Cài đặt quy trình CD (Continuous Delivery)}
Sau khi các container đơn lẻ hoàn tất, công việc PUBLISH\_TO\_GHCR sử dụng công
cụ \texttt{docker buildx imagetools} để ghép nối hai chuỗi mã băm (Digest của
AMD64 và ARM64) thành một OCI Manifest duy nhất.

Kết quả đầu ra là một Docker Image hoàn chỉnh trên kho lưu trữ GHCR. Người dùng
máy chủ truyền thống hay các dòng vi xử lý kiến trúc ARM (như Apple Silicon)
khi dùng \texttt{docker pull} ứng dụng đều sẽ nhận được container tương thích
tự động.

\ctufigure{image11.png}{1.0}{Job PUBLISH\_TO\_GHCR}{fig:job-publish}

\section{Cài đặt quy trình CD (Continuous Deployment)}
\subsection{Cài đặt máy ảo Ubuntu 24.04 bằng Vagrant}
\textbf{Vagrant được cấu hình như sau:}
\begin{itemize}
  \item \textbf{\texttt{Vagrant.configure("2")}:} Sử dụng cấu hình theo chuẩn
  phiên bản API v2 của Vagrant.
  \item \textbf{\texttt{config.vm.box = "bento/ubuntu-24.04"}:} Sử dụng hệ điều
  hành Ubuntu 24.04 LTS được đóng gói bởi Bento.
  \item \textbf{\texttt{virtualbox.memory = 4096}:} Cấp phát 4096 MB (4 GB) bộ
  nhớ RAM cho máy ảo để đáp ứng nhu cầu build image và chạy nhiều container
  cùng lúc.
  \item \textbf{\texttt{virtualbox.cpus = 4}:} Phân bổ 4 nhân CPU ảo (vCPU) giúp
  tăng tốc độ build và xử lý tác vụ của ứng dụng.
  \item \textbf{Cổng 80 (Guest) $\rightarrow$ 8080 (Host):} Cho phép truy cập
  vào Reverse Proxy (ví dụ: Nginx) từ trình duyệt máy thật thông qua địa chỉ
  \url{http://localhost:8080}.
  \item \textbf{Cổng 8080 (Guest) $\rightarrow$ 8081 (Host):} Cho phép truy cập
  trực tiếp vào dịch vụ backend/ứng dụng qua địa chỉ
  \url{http://localhost:8081}.
  \item \texttt{auto\_correct: true}: Tự động chuyển sang một cổng khả dụng khác
  trên máy Host nếu cổng được chỉ định (8080 hoặc 8081) đang bị ứng dụng khác
  chiếm dụng.
  \item \texttt{config.vm.provision}: tự động hoá cài đặt và khởi tạo dịch vụ
  Docker.
\end{itemize}

\ctufigure{image12.png}{1.0}{Cấu hình Vagrant}{fig:vagrant}

\begin{table}[!htbp]
  \centering
  \fontsize{11pt}{13.42pt}\selectfont
  \setlength{\tabcolsep}{3pt}
  \renewcommand{\arraystretch}{1.25}
  \begin{tabularx}{\textwidth}{|>{\hsize=0.8\hsize\linewidth=\hsize\raggedright\arraybackslash}X|>{\hsize=1.1\hsize\linewidth=\hsize\raggedright\arraybackslash}X|>{\hsize=1.1\hsize\linewidth=\hsize\raggedright\arraybackslash}X|}
    \hline
    \ctuTableHead{Thành phần} & \ctuTableHead{Cấu hình} & \ctuTableHead{Mục đích / Chức năng} \\
    \hline
    Virtualization Provider & VMware & Nền tảng ảo hóa máy ảo \\
    \hline
    Base OS & bento/ubuntu-24.04 (v202510.26.0) & Hệ điều hành Ubuntu 24.04 LTS \\
    \hline
    CPU / RAM & 4 vCPUs / 4096 MB RAM & Đảm bảo hiệu năng cho Docker và các dịch vụ \\
    \hline
    Port 80 (Guest) & 8080 (Host) & Truy cập Reverse Proxy (\url{http://localhost:8080}) \\
    \hline
    Port 8080 (Guest) & 8081 (Host) & Truy cập trực tiếp Backend (\url{http://localhost:8081}) \\
    \hline
    Shared Directory & /vagrant $\leftrightarrow$ Thư mục dự án & Đồng bộ mã nguồn theo thời gian thực \\
    \hline
    Container Engine & docker.io, docker-compose-v2 & Quản lý và thực thi môi trường container \\
    \hline
    Registry Authentication & GitHub Container Registry (ghcr.io) & Xác thực để kéo/đẩy container image \\
    \hline
  \end{tabularx}
  \caption{Bảng cấu hình Vagrant}
  \label{tab:vagrant}
\end{table}
\FloatBarrier

\subsection{Cài đặt docker-compose}
\subsubsection{Tổng quan kiến trúc hệ thống}
\textbf{Hệ thống được chia thành các lớp chức năng sau:}
\begin{itemize}
  \item \textbf{Lớp Web (Presentation/Proxy Layer):} vproweb (Nginx).
  \item \textbf{Lớp Ứng dụng (Application Layer):} vproapp (Java Web/Backend
  App).
  \item \textbf{Lớp Middleware \& Cache:} vpromq01 (RabbitMQ), vprocache01
  (Memcached).
  \item \textbf{Lớp Dữ liệu (Database Layer):} vprodb (MySQL Database).
  \item \textbf{Lớp Tự động hóa CD (Continuous Deployment Layer):} watchtower (Tự
  động cập nhật phiên bản mới).
\end{itemize}

\ctufigure{image13.png}{1.0}{Sơ đồ kiến trúc hệ thống Web trên nền Docker}{fig:docker-architecture}

\subsubsection{Cấu hình dịch vụ Cơ sở dữ liệu vprodb}
\textbf{Dịch vụ cơ sở dữ liệu vprodb được cấu hình như sau:}
\begin{itemize}
  \item \textbf{build.context:} Xây dựng image từ Dockerfile trong thư mục
  \nolinkurl{./Docker-files/db} (chứa các script khởi tạo schema database ban
  đầu).
  \item \textbf{ports (3306:3306):} Mở cổng chuẩn của MySQL để cho phép kết nối
  từ trong mạng container hoặc từ bên ngoài.
  \item \textbf{volumes (vprodbdata:/var/lib/mysql):} Gắn volume vprodbdata vào
  thư mục lưu trữ dữ liệu của MySQL. Đảm bảo tính toàn vẹn và bền vững của dữ
  liệu (Data Persistence), dữ liệu không bị mất khi container bị xóa hoặc khởi
  động lại.
  \item \textbf{environment:} Thiết lập mật khẩu lấy từ biến môi trường cho tài
  khoản root của cơ sở dữ liệu.
\end{itemize}

\ctufigure{image14.png}{0.83}{Cơ sở dữ liệu vprodb}{fig:vprodb}

\subsubsection{Cấu hình dịch vụ Bộ nhớ đệm vprocache01}
\textbf{Dịch vụ Bộ nhớ đệm vprocache01 được cấu hình như sau:}
\begin{itemize}
  \item \textbf{image: memcached:} Sử dụng image Memcached chính thức từ Docker
  Hub.
  \item \textbf{Mục đích:} Cung cấp cơ chế lưu trữ đệm trong RAM (In-memory
  Caching) tại cổng 11211, giúp giảm tải truy vấn cho cơ sở dữ liệu và tăng tốc
  độ phản hồi cho người dùng.
\end{itemize}

\ctufigure{image15.png}{0.40}{Bộ nhớ đệm vprocache01}{fig:vprocache}

\subsubsection{Cấu hình dịch vụ RabbitMQ}
\textbf{Dịch vụ RabbitMQ được cấu hình như sau:}
\begin{itemize}
  \item \textbf{image rabbitmq:} Sử dụng hệ thống Message Broker RabbitMQ.
  \item \textbf{ports (5672:5672):} Cổng giao thức AMQP dùng để truyền/nhận các
  bản tin (message) giữa các thành phần.
  \item \textbf{Mục đích:} Hỗ trợ xử lý bất đồng bộ (Asynchronous Processing) và
  phân tách phụ thuộc (Decoupling) giữa các module tác vụ.
\end{itemize}

\ctufigure{image16.png}{0.84}{Cấu hình dịch vụ RabbitMQ}{fig:rabbitmq}

\subsubsection{Cấu hình dịch vụ Watchtower}
Dịch vụ Watchtower được cấu hình nhằm tự động phát hiện và triển khai image mới
của ứng dụng Backend.
\begin{itemize}
  \item \textbf{image:} Sử dụng image \texttt{nickfedor/watchtower:1.21.0} để
  theo dõi phiên bản image của các container được chỉ định.
  \item \textbf{restart:} Thiết lập \texttt{unless-stopped} để dịch vụ tự khởi
  động lại khi Docker daemon hoặc máy chủ khởi động lại, trừ khi dịch vụ được
  dừng chủ động.
  \item \textbf{\texttt{/var/run/docker.sock:/var/run/docker.sock}:} Gắn Docker
  socket của máy Host để Watchtower có thể kiểm tra image, dừng và tạo lại
  container khi có bản cập nhật.
  \item \textbf{\texttt{/root/.docker/config.json:/config.json:ro}:} Gắn thông
  tin xác thực Docker ở chế độ chỉ đọc, cho phép Watchtower truy cập image
  private từ GitHub Container Registry.
  \item \textbf{\texttt{-{}-interval 60}:} Thực hiện kiểm tra registry theo chu
  kỳ 60 giây để phát hiện image tag \texttt{latest} mới.
  \item \textbf{\texttt{-{}-cleanup}:} Tự động dọn dẹp các image cũ sau khi
  container được cập nhật, giúp tiết kiệm dung lượng lưu trữ.
  \item \textbf{\texttt{-{}-debug}:} Bật log chi tiết để hỗ trợ theo dõi tiến
  trình và gỡ lỗi quá trình cập nhật.
  \item \textbf{vproapp:} Giới hạn phạm vi theo dõi cho container vproapp, tránh
  cập nhật ngoài ý muốn đối với các dịch vụ hạ tầng khác.
\end{itemize}

\ctufigure{image17.png}{0.94}{Dịch vụ Watchtower}{fig:watchtower-service}

\subsubsection{Cấu hình dịch vụ Ứng dụng Backend}
\textbf{Dịch vụ Ứng dụng Backend được cấu hình như sau:}
\begin{itemize}
  \item \textbf{image:} Kéo image ứng dụng trực tiếp từ GitHub Container
  Registry (GHCR) (\texttt{ghcr.io/duytechie/actionprofile:latest}).
  \item \textbf{labels:} Đánh dấu nhãn \texttt{enable: "true"} để dịch vụ
  Watchtower nhận diện và theo dõi việc cập nhật image mới.
  \item \textbf{depends\_on:} Chỉ định thứ tự phụ thuộc. vproapp sẽ chỉ được
  khởi động sau khi 3 dịch vụ nền tảng (vprodb, vprocache01, vpromq01) đã được
  khởi chạy.
\end{itemize}

\ctufigure{image18.png}{0.80}{Dịch vụ Ứng dụng Backend}{fig:vproapp}

\subsubsection{Cấu hình dịch vụ Web Proxy vproweb}
\textbf{Dịch vụ Web Proxy vproweb được cấu hình như sau:}
\begin{itemize}
  \item \textbf{build.context:} Xây dựng image từ thư mục
  \nolinkurl{./Docker-files/web}. Dockerfile sử dụng image Nginx làm nền tảng, xoá
  cấu hình mặc định và sao chép file \texttt{nginvproapp.conf} vào thư mục cấu
  hình của Nginx.
  \item \textbf{image:} Đặt tên image sau khi build là
  \nolinkurl{vprocontainers/vprofileweb} để thuận tiện cho việc quản lý và triển
  khai lại dịch vụ.
  \item \textbf{ports (80:80):} Công khai cổng HTTP 80 của container ra máy Host,
  tạo điểm truy cập web cho người dùng.
  \item \textbf{nginvproapp.conf:} Khai báo upstream vproapp trỏ đến
  \texttt{vproapp:8080} và sử dụng \texttt{proxy\_pass} để chuyển tiếp mọi yêu
  cầu từ Nginx đến ứng dụng Backend thông qua mạng nội bộ Docker.
  \item \textbf{depends\_on:} Chỉ định vproweb khởi động sau vproapp, bảo đảm
  lớp proxy được triển khai sau dịch vụ ứng dụng Backend.
\end{itemize}

\ctufigure{image19.png}{0.66}{Dịch vụ Web Proxy vproweb}{fig:vproweb}

\subsubsection{Theo dõi quá trình triển khai tự động}
Sau khi image ứng dụng được phát hành, quá trình triển khai có thể được theo
dõi qua kết quả workflow và log của Watchtower. Log cần thể hiện các bước phát
hiện image mới, dừng container vproapp và khởi tạo lại container với cấu hình
đã cập nhật.

\ctufigure{image20.png}{1.0}{Kết quả thực thi workflow triển khai ứng dụng}{fig:workflow-result}

\ctufigure{image21.png}{0.62}{Kết quả Workflow trong VSCode}{fig:workflow-vscode}

Log của Watchtower được sử dụng để xác nhận hệ thống đã kéo image mới và cập
nhật container vproapp thành công.

\ctufigure{image22.png}{1.0}{Log của Watchtower khi cập nhật container vproapp}{fig:watchtower-log-deploy}

\section{Cài đặt cơ chế triển khai Blue-Green}
\subsection{Tách dịch vụ ứng dụng thành Blue và Green}
\ctutodo{Đoạn \texttt{docker-compose.yml} mới; healthcheck cho container
ứng dụng.}
\subsection{Cấu hình Nginx chuyển hướng lưu lượng}
\ctutodo{Template upstream và lệnh reload.}
\subsection{Kịch bản triển khai và quay lui}
\ctutodo{Trích các phần chính của \texttt{deploy.sh}: \texttt{cosign verify},
vòng lặp health check, chuyển upstream, rollback.}
\subsection{Cài đặt self-hosted runner và job Deploy}
\ctutodo{Cài runner trên máy ảo bằng Vagrant provision; job \texttt{DEPLOY} với
\texttt{runs-on: self-hosted}, \texttt{needs: PUBLISH\_TO\_GHCR}.}

% -----------------------------------------------------------------------------
% CHƯƠNG V
% -----------------------------------------------------------------------------
\chapter{THỰC NGHIỆM VÀ ĐÁNH GIÁ}
\label{chap:evaluation}
\section{Môi trường và phương pháp thực nghiệm}
\ctutodo{Cấu hình máy host, máy ảo (4 vCPU, 4~GB RAM), runner GitHub; phiên bản
công cụ. Định nghĩa baseline (pipeline niên luận + Watchtower) và hệ thống đề
xuất. Số lần lặp mỗi thực nghiệm; cách tính trung bình và độ lệch chuẩn.}

\section{Kiểm chứng chức năng pipeline}
\subsection{Mục tiêu kiểm thử}
Kiểm thử quy trình CI/CD đã xây dựng cho ứng dụng VProfile để xác nhận một thay
đổi mã nguồn có được tích hợp, kiểm tra, đóng gói, phát hành và triển khai đúng
theo luồng tự động hay không.
\begin{itemize}
  \item \textbf{Phát hiện lỗi cấu hình:} Kiểm tra các lỗi có thể làm workflow
  không được kích hoạt, Maven không tạo được artifact, image không được phát
  hành hoặc container không được cập nhật.
  \item \textbf{Đối chiếu yêu cầu:} Đánh giá các kết quả thực tế với yêu cầu đặt
  ra cho pipeline: build một lần, quét bảo mật, build đa kiến trúc, công bố
  image lên GHCR và tự động cập nhật dịch vụ vproapp.
  \item \textbf{Quan sát log file:} Ghi nhận log Maven, kết quả GitHub Actions,
  artifact, manifest image và log Watchtower để phục vụ việc đánh giá, bảo trì
  và khắc phục sự cố về sau.
\end{itemize}

\subsection{Phạm vi kiểm thử}
Việc kiểm thử được thực hiện theo các nhóm sau:
\begin{itemize}
  \item \textbf{Kiểm thử cấu hình và kích hoạt Workflow:} Xác nhận workflow được
  kích hoạt khi có push hoặc pull request vào nhánh main, đồng thời phân quyền
  đọc mã nguồn và ghi package được khai báo đúng.
  \item \textbf{Kiểm thử tích hợp liên tục:} Xác nhận Maven có thể biên dịch,
  chạy vòng đời test, đóng gói WAR và tải artifact vprofile-app lên GitHub
  Actions.
  \item \textbf{Kiểm thử chất lượng và bảo mật:} Xác nhận báo cáo Checkstyle và
  file \texttt{trivy-results.json} được sinh ra, có thể tải lại để đánh giá mã
  nguồn và thư viện phụ thuộc.
  \item \textbf{Kiểm thử phát hành và triển khai:} Xác nhận image AMD64 và ARM64
  được gộp thành manifest trên GHCR; Watchtower phát hiện image latest mới và
  cập nhật đúng container vproapp.
\end{itemize}

Ngoài phạm vi kiểm thử:
\begin{itemize}
  \item Kiểm thử khả năng chịu tải, kiểm thử hiệu năng chuyên sâu, kiểm thử xâm
  nhập và kiểm thử đầy đủ mọi chức năng giao diện người dùng chưa thuộc phạm vi
  của niên luận này.
  \item Workflow này chỉ thực hiện các bước test cơ bản. Chưa thực hiện các quy
  trình test và security test đầy đủ và chuyên nghiệp.
\end{itemize}

\subsection{Kế hoạch kiểm thử}
Các nội dung được kiểm thử:
\begin{itemize}
  \item \textbf{Kích hoạt pipeline:} Sự kiện push và pull request vào nhánh
  main.
  \item \textbf{Build và artifact:} Lệnh \texttt{mvn install}, file WAR tại
  \texttt{target} và artifact vprofile-app.
  \item \textbf{Quét và báo cáo:} Job Security Scan, file
  \texttt{trivy-results.json} và báo cáo Checkstyle.
  \item \textbf{Image đa kiến trúc:} Job BUILD\_AMD64, BUILD\_ARM64 và
  PUBLISH\_TO\_GHCR cùng các tag commit, latest trên GHCR.
  \item \textbf{Triển khai tự động:} Chuỗi log Watchtower từ lúc phát hiện image
  mới đến khi khởi tạo lại vproapp.
\end{itemize}

\subsection{Các bước kiểm thử và kết quả ghi nhận}
Lúc này server Ubuntu để dựng trang web Java đã được khởi động sẵn.

Mã nguồn được thay đổi và đẩy lên GitHub đồng thời kích hoạt GitHub Workflow.

\ctufigure{image23.jpeg}{1.0}{Git commit}{fig:git-commit}

\subsubsection{Kiểm thử workflow và artifact}
GitHub Workflow được kích hoạt. Kết quả mong đợi là các job hiển thị Success,
\textbf{artifact vprofile-app} và \textbf{trivy-scan-results} có thể tải về, còn
manifest đa kiến trúc được tạo sau khi hai image theo digest hoàn tất.

\ctufigure{image24.png}{1.0}{Summary của GitHub Actions và danh sách Artifacts}{fig:gha-summary}

\subsubsection{Kiểm thử build và test Maven}
Kịch bản chạy lệnh \texttt{mvn test} nhằm xác nhận dự án có thể biên dịch và
các test case được Maven nhận diện. Kết quả ghi nhận BUILD SUCCESS nhưng
Tests run: 0, Failures: 0, Errors: 0.

\ctufigure{image25.png}{1.0}{Kết quả mvn test}{fig:mvn-test}

\subsubsection{Kiểm thử phát hành image đa kiến trúc}
Kịch bản kiểm tra package trên GHCR xác nhận cùng một tag \texttt{latest} có
thể tham chiếu đến image \texttt{linux/amd64} và \texttt{linux/arm64}.

\ctufigure{image26.png}{1.0}{Package trên GHCR}{fig:ghcr-package}

\subsubsection{Kiểm thử cập nhật tự động bằng Watchtower}
Sau đó nếu workflow chạy thành công thì image latest mới đã xuất hiện trên
GHCR. Kết quả mong đợi là Watchtower lần lượt kéo image, phát hiện image mới,
dừng vproapp cũ và khởi động lại vproapp với cấu hình được giữ nguyên. Các dịch
vụ hạ tầng khác không nằm trong phạm vi cập nhật của Watchtower.

Thực hiện kết nối SSH vào server Ubuntu.

\ctufigure{image27.png}{0.93}{vagrant ssh}{fig:vagrant-ssh}

Kiểm tra log của watchtower bằng lệnh sau:

\begin{lstlisting}[language=bash]
docker compose logs -f watchtower | rg --line-buffered 'Image pull completed|Found new image|Stopping container|Starting container with updated configuration|Update session completed'
\end{lstlisting}

\ctufigure{image22.png}{1.0}{Log Watchtower}{fig:watchtower-log-test}

Sau đó người dùng sẽ nhận được bản cập nhật mới một cách tự động sau lần
refresh lại trang web.

\ctufigure{image28.png}{0.56}{Phiên bản web mới được cập nhật}{fig:web-updated}

\subsubsection{Đánh giá kiểm thử}
Quy trình đã đáp ứng được mục tiêu tự động hóa việc build, test, security test,
phát hành image đa kiến trúc và triển khai lại dịch vụ của ứng dụng web.

\section{Thực nghiệm E1: Hiệu quả của các cổng bảo mật}
\ctutodo{Tạo các nhánh chứa lỗi cố ý: (a) thư viện có CVE đã biết, ví dụ
log4j-core 2.14.1; (b) secret giả; (c) đoạn mã SQL injection; (d) base image cũ
có CVE. Bảng: kịch bản $\times$ gate, baseline có chặn không / hệ thống mới có
chặn không, công cụ phát hiện, số false positive.}

\section{Thực nghiệm E2: Chi phí thời gian của pipeline}
\ctutodo{Chạy baseline và pipeline mới mỗi loại 15--20 lần; bảng và biểu đồ thời
gian từng job (trung bình $\pm$ độ lệch chuẩn); phần trăm thời gian tăng thêm;
ảnh hưởng của cache.}

\section{Thực nghiệm E3: Thời gian gián đoạn khi triển khai}
\ctutodo{k6 tạo tải liên tục (ví dụ 50 VUs) trong lúc triển khai; mỗi chiến lược
10 lần. Bảng: số request lỗi, thời gian gián đoạn (giây), độ trễ p95; biểu đồ
tỉ lệ lỗi theo thời gian.}

\section{Thực nghiệm E4: Khả năng tự phục hồi}
\ctutodo{(a) triển khai image lỗi khởi động; (b) image chưa ký. So sánh Watchtower
và Blue-Green: phiên bản lỗi có nhận lưu lượng không, thời gian khôi phục
(MTTR), image chưa ký có bị từ chối không.}

\section{Thảo luận}
\subsection{Trả lời câu hỏi nghiên cứu}
\ctutodo{Trả lời RQ1 và RQ2 dựa trên số liệu E1--E4.}
\subsection{Các yếu tố ảnh hưởng đến tính hợp lệ}
\ctutodo{Threats to validity: chỉ một ứng dụng mẫu, một máy chủ, runner GitHub có
thời gian chạy dao động, lỗi được tạo nhân tạo.}

% -----------------------------------------------------------------------------
% KẾT LUẬN
% -----------------------------------------------------------------------------
\ctuheading{KẾT LUẬN VÀ HƯỚNG PHÁT TRIỂN}
\setcounter{introsection}{0}
\introsection{Kết quả đạt được}
\ctutodo{Tóm tắt đóng góp và kết quả chính bằng số liệu.}
\introsection{Hạn chế}
\ctutodo{Một máy chủ đơn; cơ sở dữ liệu dùng chung, chưa xử lý migration; chưa
đánh giá DAST; v.v.}
\introsection{Hướng phát triển}
\ctutodo{Chuyển lên Kubernetes với Argo CD (GitOps) và Argo Rollouts (Canary);
bổ sung DAST (OWASP ZAP); chính sách admission (Kyverno) xác minh chữ ký; đạt
SLSA cấp cao hơn.}

\clearpage
\phantomsection
\addcontentsline{toc}{chapter}{TÀI LIỆU THAM KHẢO}
\begin{thebibliography}{99}
  \bibitem{vprofile}
  vprofile-project, \url{https://github.com/hkhcoder/vprofile-project}.

  \bibitem{200lab-cicd}
  200Lab, ``CI/CD là gì?'', \url{https://200lab.io/blog/ci-cd-la-gi}.

  \bibitem{push-ghcr}
  GitHub Marketplace, ``Push to GHCR'',
  \url{https://github.com/marketplace/actions/push-to-ghcr}.

  \bibitem{login-action}
  Docker, ``Login Action'', \url{https://github.com/docker/login-action}.

  \bibitem{blackduck-cicd}
  Black Duck, ``Continuous Integration and Continuous Delivery/Deployment'',

  \bibitem{nist-ssdf}
  M. Souppaya, K. Scarfone, D. Dodson, ``Secure Software Development Framework
  (SSDF) Version 1.1'', NIST SP 800-218, 2022.

  \bibitem{slsa}
  OpenSSF, ``SLSA specification v1.0'', \url{https://slsa.dev/spec/v1.0/}.

  \bibitem{owasp-cicd}
  OWASP, ``Top 10 CI/CD Security Risks'',
  \url{https://owasp.org/www-project-top-10-ci-cd-security-risks/}.

  \bibitem{dora}
  Google Cloud, ``Accelerate State of DevOps Report'', \url{https://dora.dev/research/}.

  \bibitem{fowler-bluegreen}
  M. Fowler, ``BlueGreenDeployment'',
  \url{https://martinfowler.com/bliki/BlueGreenDeployment.html}.

  \bibitem{sigstore}
  Sigstore, ``Cosign documentation'', \url{https://docs.sigstore.dev/}.

  \bibitem{trivy}
  Aqua Security, ``Trivy documentation'', \url{https://trivy.dev/}.

  \bibitem{k6}
  Grafana Labs, ``k6 documentation'', \url{https://grafana.com/docs/k6/}.

  \bibitem{mohan2016}
  V. Mohan, L. Ben Othmane, ``SecDevOps: Is It a Marketing Buzzword? Mapping
  Research on Security in DevOps'', in \textit{Proc. 11th Int. Conf. on
  Availability, Reliability and Security (ARES)}, IEEE, 2016.

  \bibitem{myrbakken2017}
  H. Myrbakken, R. Colomo-Palacios, ``DevSecOps: A Multivocal Literature
  Review'', in \textit{Software Process Improvement and Capability
  Determination (SPICE 2017)}, CCIS vol.~770, Springer, 2017.

  \bibitem{rajapakse2022}
  R. N. Rajapakse, M. Zahedi, M. A. Babar, H. Shen, ``Challenges and solutions
  when adopting DevSecOps: A systematic review'', \textit{Information and
  Software Technology}, vol.~141, 106700, 2022.

  \bibitem{enck2022}
  W. Enck, L. Williams, ``Top Five Challenges in Software Supply Chain Security:
  Observations From 30 Industry and Government Organizations'', \textit{IEEE
  Security \& Privacy}, vol.~20, no.~2, 2022.

  \bibitem{ohm2020}
  M. Ohm, H. Plate, A. Sykosch, M. Meier, ``Backstabber's Knife Collection: A
  Review of Open Source Software Supply Chain Attacks'', in \textit{Detection
  of Intrusions and Malware, and Vulnerability Assessment (DIMVA)}, LNCS,
  Springer, 2020.

  \bibitem{ladisa2023}
  P. Ladisa, H. Plate, M. Martinez, O. Barais, ``SoK: Taxonomy of Attacks on
  Open-Source Software Supply Chains'', in \textit{IEEE Symposium on Security
  and Privacy (SP)}, 2023.

  \bibitem{sigstore-paper}
  Z. Newman, J. S. Meyers, S. Torres-Arias, ``Sigstore: Software Signing for
  Everybody'', in \textit{Proc. ACM SIGSAC Conf. on Computer and Communications
  Security (CCS)}, 2022.

  \bibitem{shahin2017}
  M. Shahin, M. Ali Babar, L. Zhu, ``Continuous Integration, Delivery and
  Deployment: A Systematic Review on Approaches, Tools, Challenges and
  Practices'', \textit{IEEE Access}, vol.~5, pp.~3909--3943, 2017.

  \bibitem{humble2010}
  J. Humble, D. Farley, \textit{Continuous Delivery: Reliable Software Releases
  through Build, Test, and Deployment Automation}, Addison-Wesley, 2010.

  \bibitem{forsgren2018}
  N. Forsgren, J. Humble, G. Kim, \textit{Accelerate: The Science of Lean
  Software and DevOps}, IT Revolution Press, 2018.

  \bibitem{mandiant-sunburst}
  Mandiant (FireEye), ``Highly Evasive Attacker Leverages SolarWinds Supply
  Chain to Compromise Multiple Global Victims With SUNBURST Backdoor'', 2020.

  \bibitem{nvd-log4shell}
  NIST National Vulnerability Database, ``CVE-2021-44228'',
  \url{https://nvd.nist.gov/vuln/detail/CVE-2021-44228}.

  \bibitem{nvd-xz}
  NIST National Vulnerability Database, ``CVE-2024-3094'',
  \url{https://nvd.nist.gov/vuln/detail/CVE-2024-3094}.

  \bibitem{cvss}
  FIRST, ``Common Vulnerability Scoring System v3.1: Specification Document'',
  \url{https://www.first.org/cvss/v3.1/specification-document}.

  \bibitem{codeql}
  GitHub, ``CodeQL documentation'', \url{https://codeql.github.com/docs/}.

  \bibitem{findsecbugs}
  Find Security Bugs, ``The SpotBugs plugin for security audits of Java web
  applications'', \url{https://find-sec-bugs.github.io/}.

  \bibitem{dependency-check}
  OWASP, ``Dependency-Check'',
  \url{https://owasp.org/www-project-dependency-check/}.

  \bibitem{gitleaks}
  Gitleaks, \url{https://github.com/gitleaks/gitleaks}.

  \bibitem{spdx}
  ISO/IEC 5962:2021, ``Information technology -- SPDX\textsuperscript{\textregistered}
  Specification V2.2.1'', 2021.

  \bibitem{cyclonedx}
  OWASP, ``CycloneDX Bill of Materials Standard'', \url{https://cyclonedx.org/}.

  \bibitem{nginx-control}
  NGINX, ``Controlling nginx'', \url{https://nginx.org/en/docs/control.html}.

  \bibitem{docker-healthcheck}
  Docker, ``Dockerfile reference -- HEALTHCHECK'',
  \url{https://docs.docker.com/reference/dockerfile/\#healthcheck}.

  \bibitem{eo14028}
  The White House, ``Executive Order 14028: Improving the Nation's
  Cybersecurity'', 2021.
\end{thebibliography}

\end{document}
